\documentclass[11pt]{article}

\usepackage[margin=1.05in]{geometry}
\usepackage{amsmath,amssymb,amsthm,mathtools}
\usepackage{stmaryrd}
\usepackage{mathpartir}
\usepackage{enumitem}
\usepackage{verbatim}
\usepackage{hyperref}
\usepackage{booktabs,tabularx}
\usepackage{caption}

\hypersetup{
  colorlinks=true,
  linkcolor=blue,
  citecolor=blue,
  urlcolor=blue
}

% -----------------------------------------------------------------------------
% Macros
% -----------------------------------------------------------------------------

\newcommand{\Proc}{\mathrm{Pr}}
\newcommand{\Rew}{\mathrm{R}}
\newcommand{\Nm}{\mathrm{Nm}}
\newcommand{\ctx}{\;|\;}
\newcommand{\rw}{\rightsquigarrow}
\newcommand{\ty}{\mathrel{::}}
\newcommand{\dep}{\mathsf{\Pi}}
\newcommand{\app}{\mathsf{app}}
\newcommand{\lam}{\mathsf{lam}}
\newcommand{\Lam}{\mathsf{lam}^{\sharp}}
\newcommand{\out}{\mathsf{out}}
\newcommand{\inp}{\mathsf{in}}
\newcommand{\Inp}{\mathsf{in}^{\sharp}}
\newcommand{\parop}{\mathsf{par}}
\newcommand{\comm}{\mathsf{comm}}
\newcommand{\erase}{\mathsf{erase}}
\newcommand{\Dia}{\Diamond}
\newcommand{\hole}{[-]}
\newcommand{\Sites}{\mathsf{Sites}}
\newcommand{\Occ}{\mathsf{Occ}}
\newcommand{\Prfl}{\mathsf{Prfl}}
\newcommand{\RawPrfl}{\mathsf{RawPrfl}}
\newcommand{\TOcc}{\mathsf{TOcc}}
\newcommand{\Match}{\mathsf{Match}}
\newcommand{\Con}{\mathsf{Con}}
\newcommand{\Sig}{\Sigma}
\newcommand{\Hyp}{\mathsf{Hyp}}
\newcommand{\Ev}{\mathsf{Ev}}
\newcommand{\tm}{\mathsf{tm}}
\newcommand{\tp}{\mathsf{tp}}
\newcommand{\FV}{\mathsf{FV}}
\newcommand{\Sub}{\mathsf{Sub}_{\Proc}}
\newcommand{\Mod}[4]{\left\langle #1\right\rangle^{#2}_{#3} #4}
\newcommand{\judg}[3]{#1\ctx #2\;\vdash\;#3}

\newtheorem{definition}{Definition}[section]
\newtheorem{theorem}[definition]{Theorem}
\newtheorem{proposition}[definition]{Proposition}
\newtheorem{remark}[definition]{Remark}
\newtheorem{example}[definition]{Example}

\title{Deriving Hypercubes of Type Systems}
\author{Michael Stay\\\texttt{director.research@f1r3fly.io}
  \and L. Gregory Meredith\\\texttt{f1r3flyceo@gmail.com}
  \and Christian Wells\\\texttt{cb@holdea.net}}
\date{}

\begin{document}

\maketitle

\begin{abstract}
Barendregt’s Lambda Cube organizes typed lambda calculi by enabling different dependencies between terms and types. We show how to derive an analogous hypercube of type systems from an arbitrary finite presentation of an operational theory and a few design choices.
\end{abstract}

\section{Introduction}
\label{sec:introduction}

Suppose we start with the untyped $\lambda$-calculus.  How can we equip it with a type system?  Barendregt \cite{Barendregt1991} gave eight different type systems arranged in a cube.  The type systems adjoin to the untyped calculus two special symbols $*$ and $\square$ called ``sorts'' for tracking the difference between terms and types; they also adjoin a dependent product type code $\Pi_{x::A} B.$  All the type systems share the same axioms and inference rules; they differ only in the choice of which sorts can be used in the formation and introduction rules for the dependent product.

Suppose instead we start with the asynchronous $\pi$-calculus \cite{MilnerParrowWalker1992, HondaTokoro1991, Boudol1992}.  Is there a similar way to equip it with a set of type systems?  If so, is there a construction that can play a role similar to that of the dependent product? This paper answers in the affirmative: given a finite operational presentation of a small-step semantics, together with a finite declaration of the displayed source sites to type and a choice of allowed local sort profiles, there is a canonical completion which freely adjoins modal dependent type formers analogous to the dependent product.

This emphasis on contexts is related to the use of contextual labels in the reactive-systems approach of Leifer and Milner \cite{LeiferMilner2000}.  There, labels record the minimal contexts which enable a reaction, and this contextual presentation is used to prove that the induced bisimilarity is a congruence.  The present paper uses a similar operational insight in a different direction.  Instead of using contexts only as labels for transitions, we use displayed redex contexts as generators for modal typing structure: the context which enables a computation also determines the shape of the corresponding contextual may-type.  In this sense the construction may be viewed as a type-theoretic, or Hennessy--Milner-style, refinement of the contextual-labeling idea.

\section{Barendregt's lambda cube}
\label{subsec:barendregt-lambda-cube}

Barendregt's lambda cube presents eight typed lambda calculi by varying which dependencies are allowed among terms and types.  There are two sorts,
\[
  S = \{\ast,\square\},
\]
where \(\ast\) is the sort of ordinary types and \(\square\) is the sort of kinds.  The sort axiom common to all systems in the cube is
\begin{mathpar}
\inferrule*[right=Axiom.]
  { }
  { \vdash \ast : \square }
\end{mathpar}

Let \(V\) be a set of variables.  Raw terms are generated by
\[
  A,B,C,D
  \;::=\;
  x
  \mid
  s
  \mid
  A\,B
  \mid
  \lambda x:A.\,B
  \mid
  \Pi_{x:A}\,B ,
\]
where \(x\in V\) and \(s\in S\).  Contexts are generated by
\[
  \Gamma
  \;::=\;
  \emptyset
  \mid
  \Gamma,x:A .
\]

The one-step beta reduction relation is generated by the beta contraction rule
\begin{mathpar}
\inferrule*[right=\(\beta\)]
  { }
  { (\lambda x:A.\,B)\,C \to_{\beta} B[C/x] }
\end{mathpar}
and by compatibility with the term formers:
\begin{mathpar}
\inferrule*
  { A \to_{\beta} A' }
  { A\,B \to_{\beta} A'\,B }

\inferrule*
  { B \to_{\beta} B' }
  { A\,B \to_{\beta} A\,B' }

\inferrule*
  { A \to_{\beta} A' }
  { \lambda x:A.\,B \to_{\beta} \lambda x:A'.\,B }

\inferrule*
  { B \to_{\beta} B' }
  { \lambda x:A.\,B \to_{\beta} \lambda x:A.\,B' }

\inferrule*
  { A \to_{\beta} A' }
  { \Pi_{x:A}\,B \to_{\beta} \Pi_{x:A'}\,B }

\inferrule*[right=.]
  { B \to_{\beta} B' }
  { \Pi_{x:A}\,B \to_{\beta} \Pi_{x:A}\,B' }
\end{mathpar}
Write \(=_{\beta}\) for the beta-convertibility relation generated by
\(\to_{\beta}\).

The following typing rules are common to all eight systems:
\begin{mathpar}
\inferrule*[right=Start]
  { \Gamma \vdash A : s }
  { \Gamma,x:A \vdash x : A }
  \quad x\notin\Gamma

\inferrule*[right=Weakening]
  { \Gamma \vdash A : B \\ \Gamma \vdash C : s }
  { \Gamma,x:C \vdash A : B }
  \quad x\notin\Gamma

\inferrule*[right=Application]
  { \Gamma \vdash C : \Pi_{x:A}\,B \\ \Gamma \vdash D : A }
  { \Gamma \vdash C\,D : B[D/x] }

\inferrule*[right=Conversion.]
  { \Gamma \vdash A : B \\ B =_{\beta} B' \\ \Gamma \vdash B' : s }
  { \Gamma \vdash A : B' }
\end{mathpar}

The systems differ only in the pairs of sorts \((s_1,s_2)\) allowed in the product and abstraction rules:
\begin{mathpar}
\inferrule*[right=Product]
  { \Gamma \vdash A : s_1 \\ \Gamma,x:A \vdash B : s_2 }
  { \Gamma \vdash \Pi_{x:A}\,B : s_2 }
\\
\inferrule*[right=Abstraction.]
  { \Gamma \vdash A : s_1
    \\ \Gamma,x:A \vdash B : s_2
    \\ \Gamma,x:A \vdash C : B }
  { \Gamma \vdash \lambda x:A.\,C : \Pi_{x:A}\,B }
\end{mathpar}
When $B$ does not depend on $x,$ we write $A \to B$ instead of $\Pi_{x:A}\,B.$

The four possible pairs have the following meanings:
\[
\begin{array}{ccl}
  (\ast,\ast) &:& \text{terms may depend on terms},\\
  (\ast,\square) &:& \text{types may depend on terms},\\
  (\square,\ast) &:& \text{terms may depend on types},\\
  (\square,\square) &:& \text{types may depend on types}.
\end{array}
\]

The eight vertices of the cube are obtained by choosing which of the three non-base pairs are present, always including the simply typed base pair
\((\ast,\ast)\):
\[
\begin{array}{c|cccc}
  \text{system}
    & (\ast,\ast)
    & (\ast,\square)
    & (\square,\ast)
    & (\square,\square)
  \\ \hline
  \lambda\to
    & \times &  &  & 
  \\
  \lambda P
    & \times & \times &  & 
  \\
  \lambda 2
    & \times &  & \times & 
  \\
  \lambda\underline{\omega}
    & \times &  &  & \times
  \\
  \lambda P2
    & \times & \times & \times & 
  \\
  \lambda P\underline{\omega}
    & \times & \times &  & \times
  \\
  \lambda\omega
    & \times &  & \times & \times
  \\
  \lambda C
    & \times & \times & \times & \times
\end{array}
\]

So \(\lambda\to\) is the simply typed lambda calculus, \(\lambda P\) adds dependent types, \(\lambda 2\) adds polymorphism, \(\lambda\underline{\omega}\) adds type operators, and \(\lambda C\) (the calculus of constructions) contains all four product profiles.

\section{Operational semantics}

In order for us to be able to derive a type system for a programming language or virtual machine, we need some formal way of describing how it works.  What we describe below is morally the same as Plotkin's structured operational semantics \cite{PlotkinSOS}, but cast in the language of category theory.  

Hamilton and De Morgan introduced the notion of universal algebra \cite[preface]{Whitehead1898}, the study of sets equipped with functions satisfying equations. Graphs, monoids, groups, rings, and modules over a fixed ring are a few examples of algebraic structures.  Fields are not algebraic in this purely equational sense, because the usual definition involves the inequality requirement $0 \ne 1$ and a partial inverse operation.

In 1963, Lawvere \cite{Lawvere1963thesis} showed how the definition of an algebraic structure could be encoded into what is now called a ``Lawvere theory''.  Let $A$ (for ``arity'') be the opposite of the category of finite sets.  A {\em Lawvere theory} is a category $T$ equipped with an identity-on-objects functor $A \to T:$  an $n$-element finite set is the coproduct of $n$ singletons, so in the opposite category, it's the product of $n$ copies of a generating object.  Instances of the algebraic structure correspond to product-preserving functors from $T$ to the category of sets.  Homomorphisms of the algebraic structures correspond to natural transformations between the functors.  

For example, the theory of a monoid Th(Mon) can be presented with the following data:
\begin{itemize}
    \item a generating object $M$
    \item generating morphisms
    \begin{itemize}
        \item $e:1 \to M$
        \item $m:M^2 \to M$
    \end{itemize}
    \item equations
    \begin{itemize}
        \item $m \circ (e \times M) = M$ (left unit)
        \item $m \circ (M \times e) = M$ (right unit)
        \item $m \circ (m \times M) = m \circ (M \times m)$ (associativity)
    \end{itemize}
\end{itemize}
A product-preserving functor $F$ from Th(Mon) to Set picks out a set of elements $F(M)$ and equips it with functions $F(e):1 \to F(M)$ and $F(m):F(M)^2 \to F(M)$ that pick out the identity element and multiply two elements, respectively---that is, $F(M)$ is equipped with the structure of a monoid.  Note that because the functor is product-preserving, $F(1) \cong 1$ and $F(M^2) \cong F(M)^2.$  Functoriality preserves equations, so the multiplication function is associative and unital.  A natural transformation $h: F \Rightarrow F'$ assigns to the generating object $M$ a function $h_M: F(M) \to F'(M)$ between the two sets of elements such that the unit laws and associativity are preserved---that is, $h$ defines a monoid homomorphism.  So Th(Mon) captures ``what it means to be a monoid''.  We can use a similar approach to capture ``what it means to implement a virtual machine''.

While ordinary Lawvere theories are generated by a single object, there is a notion of ``multisorted'' Lawvere theories \cite{TrimbleMultisorted} that are generated by a finite set of ``sorts'', which here means ``generating object''---though from here on, we will use the term ``carrier'' to refer to an object of a theory, since ``sort'' has the conflicting meaning used by Barendregt above.  Their essentially algebraic analogue is presented by finite-limit theories.  The extra finite-limit structure allows one to specify definable subobjects, domains of partial operations, and evidence objects, not merely total operations.  This approach has proven very fruitful \cite{AdamekRosickyVitale}.

A graph-structured theory is one equipped with two distinguished objects \(V\) and \(E\), together with two morphisms
\[
s,t:E\to V.
\]
A model of this part of the theory contains a directed multigraph: a set of vertices, a set of edges, and source and target functions.  Parallel edges and self-loops are allowed.

\subsection{Operational theories}

A graph-structured theory is a useful setting for expressing small-step operational semantics, where we have a graph of states/terms with state updates/rewrites between them.  We will use a graph-structured, finitely complete cartesian closed setting.  The finite limits provide the structural and evidence layer, while the cartesian closed structure provides variable binding and capture-avoiding substitution.  We write \(\Proc\) for the vertex object and \(\Rew\) for the edge object.  Elements of \(\Proc\) are processes, terms, or machine states; elements of \(\Rew\) are single-step rewrites.  The maps
\[
s,t:\Rew\to\Proc
\]
assign to each rewrite its source and target process.

Let \(A\) be the free finitely complete cartesian closed category on a chosen set of carriers.  Its objects are the formal expressions freely generated from those generators using finite-limit structure and exponentials.  An operational theory is a finitely complete cartesian closed category \(T\), equipped with an identity-on-objects functor \(A\to T\), together with two distinguished carriers \(\Proc\) and \(\Rew\) and morphisms
\[
s,t:\Rew\to\Proc.
\]
The identity-on-objects condition means that \(T\) extends the available
operations, equations, and rewrite constructors while keeping the same
object expressions as the base theory \(A\).


\begin{definition}[Finite presentation of an operational theory]
A finite presentation \(\mathcal{P}\) of an operational theory consists of:
\begin{enumerate}[label=\textup{(\roman*)}]
    \item finitely many carrier objects other than \(\Proc\) and \(\Rew\),
    \item finitely many term constructor symbols between finite-limit and exponential expressions in those carriers other than $s,t:\Rew \to \Proc$,
    \item finitely many relation symbols,
    \item finitely many displayed equation schemas between compositions of the term constructors,
    \item and finitely many primitive rewrite constructors
    \[
        \Gamma_i,\Delta_i\vdash \rho_i:L_i\rw R_i.
    \]
\end{enumerate}
The displayed source term \(L_i\) and target term \(R_i\) are part of the presentation. Subterm occurrences of \(L_i\) are presentation data, not invariant structure of an abstract quotient theory.

We keep three layers of presentation data separate.  The \emph{display layer} consists of the raw syntax trees of the constructors, equations, and primitive rewrite sources and targets.  This is the layer from which source occurrences are named.  The \emph{equational layer} consists of the listed equations; it generates the ordinary equality or structural congruence of processes in models of the operational theory.  The \emph{profile-matching layer} consists of finite relations between displayed type-bearing occurrences of equations, used only by the profile algorithm below.  A displayed equation schema may include such a finite list of matched type-bearing occurrence pairs.  Absent such a list, the algorithm uses the syntactic default matching defined in Section~\ref{sec:construction}.
\end{definition}

This presentation-dependence is intentional. We will be constructing a set of modal operators from the displayed sources of rewrite rules by identifying specified source occurrences. This requires choosing a specific representative of an equivalence class for each source and choosing which displayed occurrences are operational sites. Structural equations may identify two displayed processes later, but they do not by themselves rename sites, add sites, or identify profile slots.

\subsection{The untyped \texorpdfstring{$\lambda$}{λ}-calculus}
The trivial operational theory is a theory of the simply-typed $\lambda$-calculus with the single generating type $\Proc.$  It automatically has lambda terms as part of the syntax, and beta reduction is an equality.  However, it does not have any rewrites, so it does not give us a place to stand to reason about the small-step operational semantics of $\lambda$-calculus.

To model the untyped lambda calculus, we ``wrap'' the built-in syntax with generating term constructors.  Notice below that there is no term constructor for variables, let alone a definition of free and bound variables or capture-avoiding substitution; this is all taken care of by the built-in binder and evaluation syntax.  This wrapping layer gives us the ability to be much more specific about computational aspects: we explicitly track $\beta$ reduction as a rewrite step and choose to use normal-order reduction (only allow reductions in the head position of an application) for our evaluation strategy.

Here's the presentation of the theory:

\begin{itemize}
    \item Carriers other than \(\Proc\) and \(\Rew\): None.

    \item Term constructors other than $s,t:$
    \begin{align*}
        p, q:\Proc &\vdash {\rm app}(p, q): \Proc & \mbox{application}\\
        k:\Proc \to \Proc & \vdash {\rm lam}(k): \Proc & \mbox{abstraction}\\ 
    \end{align*}
    \item Rewrite rules:
    \begin{align*}
        p:\Proc, k:\Proc \to \Proc &\vdash {\rm beta}:{\rm app}({\rm lam}(k), p) \rw k(p) & \beta \\
        q: \Proc, r: \Rew & \vdash {\rm head}: {\rm app}(s(r), q) \rw {\rm app}(t(r), q) & \mbox{head}
    \end{align*}
\end{itemize}

For the lambda calculus, Barendregt's cube may be viewed retrospectively as adjoining sorts, a dependent product type former, and inference rules for inductively defining typing.  We would like to come up with an algorithm that does most of that work for any operational theory, not just for $\lambda$-calculus.

\subsection{Asynchronous \texorpdfstring{$\pi$}{π}-calculus}
As another example of a nontrivial operational theory, we present the asynchronous \(\pi\)-calculus.  This is a simple model of concurrent execution via message passing; in it, one can represent race conditions and deadlocks.  In general, a term may reduce in many inequivalent ways: in a race, it matters who wins.

\begin{itemize}
    \item Carriers other than \(\Proc\) and \(\Rew\):
    \begin{align*}
        & \Nm & \mbox{names}
    \end{align*}
    \item Term constructors other than $s,t:$
    \begin{align*}
        &\vdash {\rm nil}: \Proc & \mbox{do-nothing} \\
        p, q: \Proc & \vdash p \;|\; q: \Proc & \mbox{parallel composition}\\
        n, m: \Nm & \vdash n!m: \Proc & \mbox{send} \\
        n: \Nm, p: \Nm \to \Proc & \vdash n?p : \Proc & \mbox{receive} \\
        p: \Proc & \vdash \;!p: \Proc & \mbox{replication}\\
        p: \Nm \to \Proc & \vdash \nu.p: \Proc & \mbox{restriction}
    \end{align*}
    
    \item Equations:
    \begin{align*}
        p: \Proc & \vdash p \;|\; {\rm nil} = p & \mbox{unit}\\
        p, q: \Proc & \vdash p \;|\; q = q \;|\; p & \mbox{commutativity}\\
        p, q, r: \Proc & \vdash (p \;|\; q) \;|\; r = p \;|\; (q \;|\; r) & \mbox{associativity}\\
        &\vdash \nu.\lambda x.{\rm nil} = {\rm nil} & \mbox{new nil}\\
        p:\Proc & \vdash \;!p = p\;|\;!p & \mbox{spawn}\\
        p:\Nm \to \Nm \to \Proc & \vdash \nu.\lambda x.\nu.\lambda y.p(x)(y) = \nu.\lambda x.\nu.\lambda y.p(y)(x) & \mbox{swap}\\
        p:\Nm \to \Proc, q: \Proc & \vdash (\nu.p)\;|\;q = \nu.\lambda x.(p(x) \;|\; q) & \mbox{scope extrusion}
    \end{align*}
    
    \item Rewrite constructors:
    \begin{align*}
        q: \Proc, r: \Rew & \vdash {\rm par}_1: s(r) \;|\; q \rw t(r) \;|\; q & \mbox{par}_1 \\
        r_1, r_2: \Rew & \vdash {\rm par}_2: s(r_1) \;|\; s(r_2) \rw t(r_1) \;|\; t(r_2) & \mbox{par}_2\\
        n, m: \Nm, p: \Nm \to \Proc &\vdash {\rm comm}: n!m\;|\;n?p \rw p(m) & \mbox{comm} \\
        r:\Nm \to \Rew &\vdash {\rm ctx}: \nu.\lambda x.s(r(x)) \rw \nu.\lambda x.t(r(x)) & \mbox{ctx}
    \end{align*}
\end{itemize}

Some notes:
\begin{itemize}

    \item We use Milner's original factorization of receive into a synchronization on a name followed by a continuation.  The target of the comm rule evaluates the continuation \(p\) at the communicated name \(m\).  Similarly, \(\nu\) is an operation on continuations; the displayed lambda notation belongs to the specified closed structure of the presentation.

    \item As an example of a race, consider the term
    \begin{align*}
        &\nu.\lambda prize.\nu.\lambda finish.\nu.\lambda Alice.\nu.\lambda Bob.\\
        &\quad(finish!Alice \;|\; finish!Bob \;|\; finish?\lambda winner.winner!prize).
    \end{align*}
    The two processes $finish!Alice$ and $finish!Bob$ race to deliver their message to the waiting continuation. The process above reduces to either
    \[
        \nu.\lambda prize.\nu.\lambda finish.\nu.\lambda Alice.\nu.\lambda Bob.(finish!Bob \;|\; Alice!prize)    
    \]
    or
    \[
        \nu.\lambda prize.\nu.\lambda finish.\nu.\lambda Alice.\nu.\lambda Bob.(finish!Alice \;|\; Bob!prize).
    \]
    These are clearly two very different outcomes.

    \item The displayed source of the communication constructor is the ordered tree \(n!m\;|\;n?p\).  Even though commutativity later identifies this process with \(n?p\;|\;n!m\), the display layer remembers the output occurrence as the left child and the input occurrence as the right child.  The reversed representative creates no additional site unless it is separately listed as a rewrite constructor or site declaration.

    \item In the scope extrusion equation and the ctx rule, the fact that $x$ does not appear to the left of the turnstile means that $x$ is fresh relative to $p,$ $q,$ and $r.$
    
    \item If the \(\pi\)-calculus used only the \({\rm par}_1\) rule, it would have interleaving semantics; \({\rm par}_2\) is also necessary for true parallel execution.
\end{itemize}

We would like to generate from this theory a new one with sorts and a generalization of the dependent product, equipped with inference rules analogous to those in the definition of the lambda cube.


\section{Dependently-typed operational theories}
\subsection{Generalizing the dependent product}
The dependent product $\Pi_{x::A}B$ can be thought of as denoting ``those processes that when applied to a term $x$ of type $A$ form a term of type $B$''.  In an operational theory, we would write it as $\dep(A, \lambda x.B),$ a morphism with two inputs: the first for the type of the dependent variable, and the second in which $x$ is bound and may appear in the type $B.$

The first thing to notice when looking for an analogous dependent type for an arbitrary operational theory is the word ``applied''.  Application is a feature of the $\lambda$-calculus; it does not appear in the $\pi$-calculus (nor in most other calculi).  Instead, $\pi$-calculus uses a commutative juxtaposition operator.  So the first generalization we have to make is to associate a dependent type to each displayed context $K$ in which a process could be placed to cause a reduction to occur.  A {\em candidate site} is a displayed subterm occurrence in the chosen source representative of a rewrite rule that we can replace with a hole to form a term context; in this paper, candidate sites have carrier $\Proc.$  A {\em declared site} is one of the candidate occurrences that the presentation explicitly selects for modal type formation.  For example, in $\lambda$-calculus, the product-like completion declares the function-position occurrence, giving contexts like $K=\app([-], x),$ while in $\pi$-calculus, the receive completion declares the input occurrence in the displayed source \(n!m\mid n?p\), giving the displayed context \(K=n!m\mid[-]\). The commutativity equation also makes \([-]\mid n!m\) extensionally equivalent after plugging a process into the hole, but it does not create a second modal site or identify it with the displayed one.

Secondly, the word ``form'' in that characterization refers to the subject reduction theorem, the fact that in the type systems of the lambda cube, reduction is type-preserving. Requiring a subject reduction theorem of all calculi is too onerous. For example, there are very few interesting properties of a $\pi$-calculus process that are invariant across all possible reductions; we are far more interested in things like ``on which names can this process send?'', and this set changes as the process evolves.  Therefore, we would like to generalize to a modal ``possibility'' type, something like ``those processes that when applied to a term $x$ of type $A$ may reduce to a term of type $B$ in one step''.  For the $\pi$-calculus, ``those processes that when juxtaposed with a process $x$ of type $A$ may reduce to a term of type $B$ in one step'' is very different from ``those processes that when juxtaposed with a process $x$ of type $A$ form a term of type $B$''.

Traditionally, the modal type ``possibly evolves to type $X$ in one step'' is denoted $\lozenge X.$  Taken together with the notation for the dependent product that includes the variable on which the type depends, these two generalizations lead us to a notation like $\langle K\rangle_{x::A}B,$ where we've placed the context $K$ inside the diamond. This dependent type should denote ``those processes that when placed in the context $K$ with a process $x$ of type $A$ may reduce to a process of type $B$ in one step''. If the subterm at a site is placed into the resulting context, it may reduce to the target of the rewrite rule, so if we factor the source of a rewrite as $K[t, x],$ judge $x$ to be of type $A$ and the target of the rewrite to be of type $B,$ we can say both $t::\langle K\rangle_{x::A}B$ and $K[t, x]: \Dia B.$

\subsection{Typed contexts}
An operational theory axiom has one variable context and one relation context.
Dependent typing requires rules whose premises and conclusions live in different
carrier-and-evidence contexts.  For example, product formation has the form
\begin{mathpar}
  \inferrule*[right={\(\Pi\)-form.}]
    {\Gamma\mid\Delta\vdash A\ty u \\
     \Gamma,x:\Proc\mid\Delta,e_x:x\ty A\vdash B\ty v}
    {\Gamma\mid\Delta\vdash \dep(A, \lambda x.B)\ty v}
\end{mathpar}
The second premise lives in the extended context
\(\Gamma,x:\Proc\mid\Delta,e_x:x\ty A\), while the conclusion lives in
\(\Gamma\mid\Delta\).  This is not a same-context axiom; it is an inference rule
that discharges a carrier variable together with explicit typing evidence for
that variable.

Accordingly, the construction below does not return a mere operational theory. It returns a finite presentation of a \emph{dependently-typed operational theory}, a presentation of an operational theory equipped with typing judgments and specified inference rules.

\subsection{Operational sites and local cubes}
\label{sec:construction}

The examples above already contain the whole idea.  A rewrite rule does not only
say that one process may step to another; its displayed source also shows where
an interaction can happen.  The completion turns a chosen occurrence in that
source into a modal dependent type former.

For beta reduction, the displayed source is
\[
  \app(\lam(t),a)\rw t(a).
\]
If we choose the function-position occurrence \(U=\lam(t)\), the surrounding
context is
\[
  K[-]=\app([-],a).
\]
The generated type says that a process behaves like a function when, after an
argument is supplied and the process is placed in this displayed application
context, one beta step can reach the target type.  This is the application-site
origin of the dependent product.

For communication in the asynchronous \(\pi\)-calculus, the displayed source is
\[
  n!m\mid n?k \rw k(m).
\]
If we choose the input occurrence \(U=n?k\), the surrounding context is
\[
  K[-]=n!m\mid[-].
\]
The resulting receive type says that an input process, when placed next to an
output supplying a message \(m\) along channel \(n\), may communicate once and
reach the continuation type.  The continuation parameter \(k\) is not supplied
by the context; it remains an ambient parameter of the rule.

These two examples also show why sites cannot simply be ``all process
subterms.''  The lambda-cube comparison wants the application/operator site, not
the argument variable or the whole beta redex.  The receive example may choose
the input site without also choosing the output site.  Thus site selection is
part of the input data.

\begin{definition}[Declared sites]
A site-marked operational presentation consists of a finite operational
presentation \(\mathcal P\) together with a finite site specification
\(\mathfrak S\).  For each primitive rewrite constructor
\[
  \Gamma_i,\Delta_i\vdash \rho_i:L_i\rw R_i,
\]
the set \(\mathfrak S_i\) consists of selected occurrence addresses in the
\emph{displayed} source tree \(L_i\) whose subterms have carrier \(\Proc\).
For a site \(s=(i,o)\), write \(U_s\) for the selected subterm and
\(K_s[-]\) for the displayed one-hole context satisfying
\[
  L_i=K_s[U_s]
\]
as raw displayed syntax.
\end{definition}

Structural equations are deliberately kept separate from these addresses.  An
equation such as associativity or commutativity of \(\mid\) may identify terms in
models, but it does not move holes, create new sites, or merge declared sites.
Likewise, it does not identify sort-profile slots unless the presentation also
contains explicit profile-matching data.  The fully formal version of this
separation, including the default site policy and the finite matching relation,
is given in Appendix~\ref{app:site-profile-details}.

A site has three kinds of variables.  The variables shared by the selected
subterm and the context are the \emph{index variables}
\[
  \vec x_s=\FV(K_s)\cap\FV(U_s).
\]
The variables supplied only by the surrounding context are the \emph{rely
variables}
\[
  \vec y_s=\FV(K_s)\setminus\FV(U_s).
\]
The variables needed by the selected subterm or the target, but not supplied by
the one-hole context, are the \emph{ambient variables}
\[
  \vec z_s=
  (\FV(U_s)\setminus\FV(K_s))\cup(\FV(R_i)\setminus\FV(L_i)).
\]
Their declarations form an ambient telescope \(\Theta_s\).  Ambient variables
are ordinary parameters; they are not axes of the hypercube.

The type former associated to a site is written
\[
  \Mod{K_s}{\vec x\ty\vec A}{\vec y\ty\vec B}{C}.
\]
This displayed type notation records the typed interface, not a context entry.
The corresponding context always separates carriers from evidence.  If
\(\vec \aleph_s\) and \(\vec \beth_s\) are the carriers of the index and rely variables,
then the carrier-and-evidence telescope is index-first:
\[
  \Gamma,\vec x:\vec \aleph_s,\vec y:\vec \beth_s
  \mid
  \Delta,\vec e_x:\vec x\ty\vec A,\vec e_y:\vec y\ty\vec B.
\]
It should be read as follows: after the index data \(\vec x\), together with
index evidence \(\vec e_x\), are fixed, and after the rely data \(\vec y\),
together with rely evidence \(\vec e_y\), are supplied, placing the process in
the displayed context \(K_s[-]\) may take one step to a process of type \(C\).
The rely types may depend on the index carrier variables.

A declared site records the context for the site interaction itself.  It does
not, by itself, say that arbitrary rewrites can be propagated through that
context.  When the operational presentation explicitly supplies a
contextual-closure constructor
\[
  r:P\rw Q
  \quad\longmapsto\quad
  \kappa_s(r):K_s[P]\rw K_s[Q],
\]
we write
\[
  \kappa_s\Vdash K_s[-]
\]
and say that the site has \emph{closure support}.  This notation is a side
condition on the operational presentation, not a typing judgment and not a
consequence of structural congruence.  The full formal definition, including
ambient, index, and rely parameters, is given in
Definition~\ref{def:closure-support}.

\paragraph{Congruence transport for rewrites.}
Closure support is stable under equality of source and target terms.  Thus,
if \(r:P\rw Q\), and the equational part of the presentation gives
\[
  P=P'
  \qquad\text{and}\qquad
  Q=Q',
\]
then the same rewrite evidence may be transported to evidence
\[
  r':P'\rw Q'.
\]
Equivalently, rewrite evidence is interpreted modulo the structural
equations of the presentation.  This transport principle is separate from
site extraction: structural equations may transport already existing rewrite
evidence, but they do not create new displayed site addresses or force new
profile-slot identifications.

The local cube for a site records which sorts are permitted for the index types,
rely types, and target type.  Its raw slots are
\[
  I_s=
  \{i_x\mid x\in\vec x_s\}\amalg
  \{r_y\mid y\in\vec y_s\}\amalg
  \{\tau_s\}.
\]
A raw profile assigns to each slot either \(\ast\) or \(\square\), with the carrier superscript understood.  We write profiles as \((\vec\alpha,\vec\beta;\gamma)\), where \(\vec\alpha\) records the sorts of the index variables and \(\vec\beta\) records the sorts of the rely variables.  The displayed matching data of the presentation then generates a finite equivalence relation on slots; a quotient profile is a raw profile constant on those equivalence classes.  A signature choice \(\Sig_s\) is any chosen subset of the resulting quotient profile space.  A based signature choice is one containing the all-\(\ast\) element.  The lambda cube and receive hypercube examples below use based signature choices.  The complete finite graph algorithm for this quotient is in Appendix~\ref{app:site-profile-details}.

\begin{example}[A nontrivial quotient]
Suppose a site has three raw slots \(p,q,r\), but the displayed matching data
forces \(p=q\).  The raw profile cube has \(2^3=8\) vertices.  The quotient
profile space has only the four profiles satisfying \(p=q\):
\[
  (\ast,\ast;\ast),\quad
  (\ast,\ast;\square),\quad
  (\square,\square;\ast),\quad
  (\square,\square;\square).
\]
Thus quotienting is not a semantic closure over all provable equations; it is a
finite, displayed, presentation-level operation.
\end{example}

\begin{center}
\small
\begin{tabularx}{\textwidth}{@{}p{0.22\textwidth}XX@{}}
\toprule
 & Untyped \(\lambda\)-calculus completion
 & Asynchronous \(\pi\)-calculus completion \\
\midrule
Operational redex
&
\[
  \app(\lam(t),a)\rw t(a)
\]
&
\[
  n!m \mid n?k \rw k(m)
\]
\\

Declared source site
&
The function-position occurrence \(\lam(t)\) in the displayed source
\(\app(\lam(t),a)\).
&
The input occurrence \(n?k\) in the displayed source
\(n!m\mid n?k\).
\\

Displayed context
&
\[
  K[-]=\app([-],a)
\]
&
\[
  K[-]=n!m\mid[-]
\]
\\

Index variables
&
None for the product site.
&
The channel name \(n\).
\\

Rely variables
&
The argument \(a\).
&
The communicated name \(m\).
\\

Ambient variables
&
The abstraction body \(t:\Proc\to\Proc\).
&
The continuation \(k:\Nm\to\Proc\).
\\

Generated contextual type
&
\[
  \Pi_{x:A}B
\]
equivalently
\[
  \Mod{\app(\hole,x)}{}{x\ty A}{B}.
\]
&
\[
  \Mod{n!m\mid\hole}{n\ty A}{m\ty B}{C}.
\]
\\

Introduction principle
&
If, under \(x\ty A\), the body \(t(x)\) has type \(B\), then the annotated
abstraction has product type:
\[
  \Lam(A,t)\ty \Pi_{x:A}B .
\]
&
If, under \(n\ty A\) and \(m\ty B\), the continuation \(k(m)\) has type \(C\),
then the annotated input has receive type:
\[
  n?_B k
  \ty
  \Mod{n!m\mid\hole}{n\ty A}{m\ty B}{C}.
\]
\\

Elimination principle
&
Application is effectful:
\[
\begin{array}{c}
  f\ty \Pi_{x:A}B
  \qquad
  a\ty A
  \\[2pt]
  \hline
  \app(f,a)\ty \Dia(B[a/x])
\end{array}
\]
&
Communication is effectful:
\[
\begin{array}{c}
  P\ty\Mod{n!m\mid\hole}{n\ty A}{m\ty B}{C} \quad
  a\ty B
  \\[2pt]
  \hline
  n!a\mid P\ty\Dia(C[a/m])
\end{array}
\]
\\
\\

Hypercube of profiles
&
Product profiles are pairs \((u,v)\in\{\ast,\square\}^2\).  Based choices
recover the Lambda-cube product-profile geometry.
&
Receive profiles are triples
\((\alpha,\beta;\gamma)\in
\{\ast^{\Nm},\square^{\Nm}\}
\times
\{\ast^{\Nm},\square^{\Nm}\}
\times
\{\ast^{\Proc},\square^{\Proc}\}\).
Based choices give \(2^7\) receive completions.
\\

Conceptual reading
&
A product type classifies processes that, when placed in argument position,
may reduce to a term of the specified result type.
&
A receive type classifies input processes that, when placed next to a matching
output, may reduce to a continuation result of the specified type.
\\
\bottomrule
\end{tabularx}
\captionof{table}{Comparison of the lambda product-site completion and the receive-site
completion for asynchronous communication.  In both cases, a displayed
operational redex site determines a contextual may-type; the Lambda-cube
dependent product is the special case generated by the beta-redex function
position.}
\label{tab:lambda-pi-completions}
\end{center}

\section{The hypercube completion}
\label{sec:hypercube-functor}

We can now summarize the construction without the bookkeeping.  An input object
is a triple
\[
  A=(\mathcal P,\mathfrak S,\Sig),
\]
where \(\mathcal P\) is a finite operational presentation, \(\mathfrak S\) is a
finite declaration of sites, and \(\Sig\) chooses allowed quotient profiles for
those sites.  The completion adjoins sorts, proof-relevant typing evidence
objects, Church-style annotations for binding constructors, the one-step
possibility type \(\Dia\), and the contextual modalities selected by \(\Sig\).

\begin{theorem}[Hypercube completion]
For every finite site-and-signature presentation
\(A=(\mathcal P,\mathfrak S,\Sig)\), there is a generated dependently typed
operational presentation
\[
  \mathcal H(A)=\mathcal H(\mathcal P,\mathfrak S,\Sig).
\]
Its generated type formers are the one-step possibility type and, for each
declared site \(s\) with \(\Sig_s\neq\varnothing\), one contextual modality
associated to the displayed site context \(K_s[-]\).  The elements of
\(\Sig_s\) do not create additional object-language constructors; they specify
the allowed sort-profile instances of the formation, introduction,
elimination, and source-subterm rules for that site.
\end{theorem}

\begin{proof}[Construction sketch]
For each relevant carrier \(X\), adjoin the sort constants
\(\ast^X,\square^X\) and the proof-relevant typing display
\(p_X:\Ev_X\to X\times X\).  For the identity-hole context, adjoin
\(\Dia C\).  For each declared site \(s\) with \(\Sig_s\neq\varnothing\),
adjoin a contextual type former
\[
  \Mod{K_s}{\vec x\ty\vec A}{\vec y\ty\vec B}{C}.
\]
For each chosen profile \(\omega\in\Sig_s\), adjoin the corresponding
profile-indexed formation, introduction, elimination, and source-subterm rule
schemas.  Binding constructors are refined by Church-style annotations.
Delayed elimination through a context is an admissible principle, available
only when the presentation supplies explicit closure support for that context.
The formal list of generated symbols and rules is given in
Appendices~\ref{app:generated-categories}
and~\ref{app:site-profile-details}; the main rules are displayed in the next
sections.
\end{proof}

\begin{proposition}[Finiteness]
If \(\mathcal P\), \(\mathfrak S\), and \(\Sig\) are finite, then
\(\mathcal H(\mathcal P,\mathfrak S,\Sig)\) is finitely generated.
\end{proposition}

\begin{proof}[Proof sketch]
There are finitely many primitive rewrite constructors and finitely many
declared site occurrences.  Each source tree has finitely many free variables,
so each site has finitely many index, rely, and ambient variables.  The raw
profile cube for a site is finite, and the consistency quotient is computed from
a finite graph.  Since there are finitely many declared sites and each \(\Sig_s\) is finite, only
finitely many contextual modal formers and finitely many profile-indexed rule
schemas are added.
\end{proof}


\section{Typing and structural rules}

For each relevant carrier object \(X\), the generated presentation contains a
proof-relevant typing display
\[
  p_X:\Ev_X\longrightarrow X\times X,
  \qquad
  p_X=\langle\tm_X,\tp_X\rangle .
\]
The map \(p_X\) is not assumed to be monic.  Its fiber over \((a,b)\) is
the object of evidence that \(a\) has classifier \(b\).  A typing-evidence
assumption is therefore written with a name:
\[
  e:a\ty^X b.
\]
When the carrier is clear, the superscript is suppressed.  A carrier declaration
and its typing evidence are kept distinct.  Thus a context containing a typed
variable is written
\[
  \Gamma,x:X\mid\Delta,e_x:x\ty^X A.
\]
In conclusions and premises of inference rules we may still write
\(\Gamma\mid\Delta\vdash a\ty^X A\) when the name of the evidence produced by
the derivation is immaterial.  In the context itself, however, the evidence
assumption is explicit.

For every relevant carrier \(X\), the sort rule is
\begin{mathpar}
  \inferrule*[right={sort.}]
    { }
    {\Gamma\mid\Delta\vdash \ast^X\ty^X\square^X}
\end{mathpar}

\paragraph{The metavariable \(J\).}
In the structural rules below, \(J\) is a metavariable for an atomic judgment of the generated presentation.  It is not a new object-language judgment former.  It ranges over the two judgmental forms of an operational theory: term judgments and relation judgments, with equations treated as a distinguished relation judgment.  Thus, metasyntactically,
\[
  J
  \;::=\;
  t:X
  \;\mid\;
  \phi(t_1,\ldots,t_n)
  \;\mid\;
  t=t',
\]
where \(X\) ranges over carrier objects and \(\phi\) ranges over relation symbols.  Typing judgments \(M\ty A\) are displayed evidence judgments associated to the generated maps \(p_X:\Ev_X\to X\times X\).  For example, \(M\ty A\), \(p=q\), and \(r:P\rw Q\) are instances of \(J\).  The rewrite judgment \(r:P\rw Q\) is still only syntactic sugar for the three atomic judgments
\[
  r:\Rew,
  \qquad
  s(r)=P,
  \qquad
  t(r)=Q.
\]
Weakening and substitution act on \(J\) by acting on every term appearing in the judgment.  In particular, \(J[a/x]\) denotes the result of capture-avoiding substitution of \(a\) for \(x\) throughout the atomic judgment \(J\).

The structured presentation has ordinary structural behavior for typed
carrier-and-evidence contexts.
\begin{itemize}
    \item Variable:
    \begin{mathpar}
      \inferrule*[right={var.}]
        {\Gamma\mid\Delta\vdash A\ty^X u}
        {\Gamma,x:X\mid\Delta,e_x:x\ty^X A\vdash x\ty^X A}
    \end{mathpar}
    
    \item Weakening:
    \begin{mathpar}
      \inferrule*[right={weak.}]
        {\Gamma\mid\Delta\vdash J \\
         \Gamma\mid\Delta\vdash A\ty^X u}
        {\Gamma,x:X\mid\Delta,e_x:x\ty^X A\vdash J}
    \end{mathpar}

    \item Substitution:
    \begin{mathpar}
      \inferrule*[right={sub.}]
        {\Gamma,x:X\mid\Delta,e_x:x\ty^X A\vdash J \\
         \Gamma\mid\Delta\vdash a\ty^X A}
        {\Gamma\mid\Delta\vdash J[a/x]}
    \end{mathpar}
\end{itemize}

\section{Single-step rewriting and possibility}

A rewrite judgment has the form
\[
  r:P\rw Q.
\]
No reflexive-transitive closure is added by the completion. If a calculus wants reflexive, transitive, or symmetric rewrite witnesses, they must be supplied explicitly in the input presentation as primitive constructors or generated equations.  Likewise, structural equations are not treated as rewrite witnesses and do not synthesize additional contextual sites.  They remain equations, used by the ordinary conversion principle below.

The identity context \(\hole\) gives the one-step possibility type
\[
  \Dia C := \Mod{\hole}{}{}{C}.
\]
This is generated separately from the declared contextual sites.  If the whole source occurrence of a primitive rewrite is explicitly declared as a site, it produces the same identity-hole context, but the default site policy excludes whole-source occurrences to avoid duplicating \(\Dia\).

Formation:
\begin{mathpar}
  \inferrule*[right={\(\Dia\)-form.}]
    {\Gamma\mid\Delta\vdash C\ty\gamma}
    {\Gamma\mid\Delta\vdash \Dia C\ty\gamma}
\end{mathpar}

Introduction:
\begin{mathpar}
\inferrule*[right={\(\Dia\)-intro}]
  {
    \Gamma\mid\Delta\vdash r:P\rw Q
    \\
    \Gamma\mid\Delta\vdash Q\ty C
  }
  {
    \Gamma\mid\Delta\vdash P\ty\Dia C
  }
\end{mathpar}
The expression \(P\ty\Dia C\) means that \(P\) can take one rewrite step to something of type \(C\). There is no automatic return rule \(P\ty C\Rightarrow P\ty\Dia C\), unless reflexive rewrite instances are supplied in the input presentation.

Elimination:
\begin{mathpar}

\inferrule*[right={\(\Dia\)-elim}]
  {
    \Gamma\mid\Delta\vdash P\ty\Dia C
    \\
    \Gamma,Q:\Proc
    \mid
    \Delta,r:P\rw Q,e_C:Q\ty C
    \vdash J
  }
  {
    \Gamma\mid\Delta\vdash J
  }
  \quad
  (Q,r,e_C\notin\FV(J)).
\end{mathpar}
The eliminator for \(\Dia\) is the existential eliminator for one-step
reachability.  From \(P\ty\Dia C\), one may reason under fresh data
\[
  Q:\Proc,\qquad r:P\rw Q,\qquad e_C:Q\ty C .
\]
The conclusion must not depend on the chosen reduct or rewrite witness.
Thus the rule uses some witness supplied by the \(\Dia\)-evidence, but does
not require the derivation of \(P\ty\Dia C\) to have been obtained by a
particular displayed \(\Dia\)-introduction step.


\section{Context-possibility types}
\label{sec:context-possibility-types}

Let \(s\in\Sites(\mathcal{P})\) be a site. Write
\[
  L_i=K_s[U_s]
\]
for the corresponding source factorization. Let
\[
  \vec{x}=\vec{x}_s,
  \qquad
  \vec{y}=\vec{y}_s,
  \qquad
  \vec{z}=\vec{z}_s.
\]
Write \(\Theta_s\) for the corresponding ambient telescope.  The generic modal
formation, introduction, and elimination rules are stable under arbitrary outer
contexts; any ambient parameters needed by \(P\), \(Q\), or \(C\) are simply part
of \(\Gamma\).  The site-specific source-subterm rule later displays \(\Theta_s\)
explicitly, because that rule uses the distinguished subterm \(U_s\) and target
\(R_i\). The site \(s\) determines the contextual modal type former
\[
  \Mod{K_s}{\vec{x}\ty\vec{A}}{\vec{y}\ty\vec{B}}{C}.
\]
For a chosen profile class
\[
  \omega=(\vec{\alpha},\vec{\beta};\gamma)\in\Sig_s,
\]
where \(\vec\alpha\) records the sorts of the index variables and
\(\vec\beta\) records the sorts of the rely variables, the construction adds the corresponding formation, introduction, elimination, and source-subterm rule instances for this already fixed type former.

The notation for contextual modalities is analogous to the notation for
dependent products.  The generated dependent product constructor is a morphism
\[
  \Pi:\Proc\times\Proc^{\Proc}\longrightarrow\Proc,
\]
and the displayed type
\[
  \Pi_{y\ty B}C
\]
is syntactic sugar for
\[
  \Pi(B,\lambda y.C).
\]
Similarly, a modal site \(s\), with context \(K_s[-]\), generates an actual
constructor in the operational theory; the displayed notation
\[
  \langle K_s\rangle^{\vec{x}\ty\vec{A}}_{\vec{y}\ty\vec{B}}C
\]
is sugar for applying that constructor to the index types, the rely-type
family, and the target-type family.  The displayed notation names the typed
interface.  In formal rules the matching carrier-and-evidence telescope is
written explicitly:
\[
  \Gamma,\vec{x}:\vec X_s,\vec{y}:\vec Y_s
  \mid
  \Delta,\vec e_x:\vec{x}\ty\vec{A},\vec e_y:\vec{y}\ty\vec{B}.
\]
Index variables are written in the superscript and come first, while rely
variables are written in the subscript and are discharged by the modal type
former.  Since the rely types may depend on the index variables, the judgments
for \(\vec{B}\) are formed after the carrier declarations and evidence
assumptions for \(\vec{x}\), and the target type \(C\) is formed after both the
index and rely declarations.  If ambient variables are needed in a particular
instance, they are included in the outer carrier context \(\Gamma\), together with
any ambient evidence assumptions in \(\Delta\).

Let \(\vec{\aleph}\) be the product of the carriers of the index variables
\(\vec{x}\), and let \(\vec{\beth}\) be the product of the carriers of the
rely variables \(\vec{y}\).  The modal constructor associated to \(s\) has
signature
\[
  \langle K_s\rangle:
  \vec{\aleph}\times
  \vec{\beth}^{\vec{\aleph}}\times
  \Proc^{\vec{\aleph}\times\vec{\beth}}
  \longrightarrow
  \Proc^{\vec{\aleph}}.
\]
Given index types \(\vec{A}:\vec{\aleph}\), rely types depending on the index
variables,
\[
  \lambda\vec{x}.\vec{B}:\vec{\beth}^{\vec{\aleph}},
\]
and a target type family
\[
  \lambda(\vec{x},\vec{y}).C:
  \Proc^{\vec{\aleph}\times\vec{\beth}},
\]
the displayed modal type is the resulting \(\Proc\)-term in the index context,
over whatever outer context \(\Gamma\) is present:
\[
  \langle K_s\rangle^{\vec{x}\ty\vec{A}}_{\vec{y}\ty\vec{B}}C
  \;:=\;
  \langle K_s\rangle
  \bigl(
    \vec{A},
    \lambda\vec{x}.\vec{B},
    \lambda(\vec{x},\vec{y}).C
  \bigr)(\vec{x}).
\]
When there are no index variables, \(\vec{\aleph}\) is terminal and the
codomain \(\Proc^{\vec{\aleph}}\) is just \(\Proc\).  In the lambda-calculus
case, the application site has no index variables and one rely variable of
carrier \(\Proc\), so this specializes to
\[
  \langle\app(\hole,y)\rangle:
  \Proc\times\Proc^{\Proc}\longrightarrow\Proc,
\]
which is exactly the dependent-product constructor.


The index variables \(\vec{x}\) remain in the external context. The rely
variables \(\vec{y}\), and their evidence variables \(\vec e_y\), are discharged
by the modal type former:
\begin{mathpar}
  \inferrule*[right={modal-form.}]
    {\Gamma\mid\Delta\vdash \vec{A}\ty\vec{\alpha} \\
     \Gamma,\vec{x}:\vec X_s
       \mid \Delta,\vec e_x:\vec{x}\ty\vec{A}
       \vdash \vec{B}\ty\vec{\beta} \\
     \Gamma,\vec{x}:\vec X_s,\vec{y}:\vec Y_s
       \mid \Delta,\vec e_x:\vec{x}\ty\vec{A},\vec e_y:\vec{y}\ty\vec{B}
       \vdash C\ty\gamma}
    {\Gamma,\vec{x}:\vec X_s
       \mid \Delta,\vec e_x:\vec{x}\ty\vec{A}
       \vdash
     \Mod{K_s}{\vec{x}\ty\vec{A}}{\vec{y}\ty\vec{B}}{C}\ty\gamma}
\end{mathpar}
Here \(\Gamma\mid\Delta\vdash\vec{A}\ty\vec{\alpha}\) abbreviates the telescope
of judgments saying that each index type has the corresponding sort.  The
premise
\[
  \Gamma,\vec{x}:\vec X_s
  \mid \Delta,\vec e_x:\vec{x}\ty\vec{A}
  \vdash \vec{B}\ty\vec{\beta}
\]
abbreviates the corresponding telescope for the rely types in the index
carrier-and-evidence context; in particular, a rely type may mention variables
already in \(\Gamma\) and the index carrier variables, but not the discharged rely
variables themselves.

Modal introduction:
Modal introduction:
\begin{mathpar}
  \inferrule*[right={modal-intro.}]
    {
      \Gamma\mid\Delta\vdash \vec A\ty\vec\alpha
      \\
      \Gamma,\vec{x}:\vec X_s
      \mid
      \Delta,\vec e_x:\vec{x}\ty\vec A
      \vdash
      \vec B\ty\vec\beta
      \\
      \Gamma,\vec{x}:\vec X_s,\vec{y}:\vec Y_s
      \mid
      \Delta,\vec e_x:\vec{x}\ty\vec A,
             \vec e_y:\vec{y}\ty\vec B
      \vdash
      C\ty\gamma
      \\
      \Gamma,\vec{x}:\vec X_s
      \mid
      \Delta,\vec e_x:\vec{x}\ty\vec A
      \vdash
      P:\Proc
      \\
      \Gamma,\vec{x}:\vec X_s,\vec{y}:\vec Y_s
      \mid
      \Delta,\vec e_x:\vec{x}\ty\vec A,
             \vec e_y:\vec{y}\ty\vec B
      \vdash
      r:K_s[P]\rw Q
      \\
      \Gamma,\vec{x}:\vec X_s,\vec{y}:\vec Y_s
      \mid
      \Delta,\vec e_x:\vec{x}\ty\vec A,
             \vec e_y:\vec{y}\ty\vec B
      \vdash
      Q\ty C
    }
    {
      \Gamma,\vec{x}:\vec X_s
      \mid
      \Delta,\vec e_x:\vec{x}\ty\vec A
      \vdash
      P\ty\Mod{K_s}{\vec{x}\ty\vec A}{\vec{y}\ty\vec B}{C}
    }
\end{mathpar}

Modal elimination:
\begin{mathpar}
  \inferrule*[right={modal-elim.}]
    {\Gamma,\vec{x}:\vec X_s
       \mid \Delta,\vec e_x:\vec{x}\ty\vec{A}
       \vdash P\ty\Mod{K_s}{\vec{x}\ty\vec{A}}{\vec{y}\ty\vec{B}}{C} \\
     \Gamma,\vec{x}:\vec X_s
       \mid \Delta,\vec e_x:\vec{x}\ty\vec{A}
       \vdash \vec{b}\ty\vec{B}}
    {\Gamma,\vec{x}:\vec X_s
       \mid \Delta,\vec e_x:\vec{x}\ty\vec{A}
       \vdash K_s[P][\vec{b}/\vec{y}]\ty
       \Dia\bigl(C[\vec{b}/\vec{y}]\bigr)}
\end{mathpar}
In the second premise, the notation \(\vec{b}\ty\vec{B}\) means that the rule is
supplied with carrier arguments \(\vec b\) together with derivations of the
corresponding rely typings.  These premise derivations are the evidence fed to
the contextual-modality evidence.

The primitive rewrite constructor itself supplies a useful derived introduction
rule. If \(s=(i,o)\) and \(L_i=K_s[U_s]\), then, in the ambient, index, and rely
context, \(\rho_i\) gives
\[
  \Gamma,\Theta_s,\vec{x}:\vec X_s,\vec{y}:\vec Y_s
  \mid
  \Delta,\vec e_x:\vec{x}\ty\vec{A},\vec e_y:\vec{y}\ty\vec{B}
  \vdash
  \rho_i:K_s[U_s]\rw R_i.
\]
Therefore modal introduction yields
\begin{mathpar}
  \inferrule*[right={source-subterm-intro.}]
    {
     \Gamma,\Theta_s\mid\Delta\vdash \vec{A}\ty\vec{\alpha}
     \\
     \Gamma,\Theta_s,\vec{x}:\vec X_s
       \mid \Delta,\vec e_x:\vec{x}\ty\vec{A}
       \vdash \vec{B}\ty\vec{\beta}
     \\
     \Gamma,\Theta_s,\vec{x}:\vec X_s,\vec{y}:\vec Y_s
       \mid
       \Delta,\vec e_x:\vec{x}\ty\vec{A},
              \vec e_y:\vec{y}\ty\vec{B}
       \vdash C\ty\gamma
     \\
     \Gamma,\Theta_s,\vec{x}:\vec X_s,\vec{y}:\vec Y_s
       \mid
       \Delta,\vec e_x:\vec{x}\ty\vec{A},
              \vec e_y:\vec{y}\ty\vec{B}
       \vdash R_i\ty C
    }
    {
     \Gamma,\Theta_s,\vec{x}:\vec X_s
       \mid \Delta,\vec e_x:\vec{x}\ty\vec{A}
       \vdash
       U_s\ty\Mod{K_s}{\vec{x}\ty\vec{A}}{\vec{y}\ty\vec{B}}{C}
    }
\end{mathpar}
This is the rule that turns a declared source site of a rewrite into a
contextual may-precondition type.  The ambient telescope \(\Theta_s\) is essential:
it accounts for the free variables of \(U_s\) that are not supplied by
\(K_s[-]\), as well as any primitive-rule parameters needed to form the target.
Its premise uses the same index-first order as modal introduction, with the
ambient telescope placed before the index declarations, because the target type
may depend on ambient variables, shared indices, and supplied rely variables.
Ambient carrier declarations live in \(\Theta_s\); any ambient evidence
assumptions live in \(\Delta\).

\paragraph{Computation for canonical source introductions.}
When rules and their evidence are retained proof-relevantly, the generated
presentation includes the computation equation saying that source-subterm
introduction followed by modal elimination is the one-step possibility proof
obtained from the displayed primitive rewrite.  Let \(d\) be a
derivation of
\[
  \Gamma,\Theta_s,\vec{x}:\vec X_s,\vec{y}:\vec Y_s
  \mid
  \Delta,\vec e_x:\vec{x}\ty\vec{A},\vec e_y:\vec{y}\ty\vec{B}
  \vdash R_i\ty C,
\]
and let \(\operatorname{intro}^{\rho_i}_s(d)\) denote the corresponding
source-subterm introduction derivation.  If \(\vec b\) is a tuple of rely
arguments, supplied together with rely-typing evidence, then
\[
  \operatorname{elim}_s
  \bigl(\operatorname{intro}^{\rho_i}_s(d),\vec b\bigr)
  =
  \Dia\text{-}\operatorname{intro}
  \bigl(\rho_i[\vec b/\vec y],d[\vec b/\vec y]\bigr).
\]
The two sides are derivations of the same judgment
\[
  \Gamma,\Theta_s,\vec{x}:\vec X_s
  \mid
  \Delta,\vec e_x:\vec{x}\ty\vec{A}
  \vdash
  K_s[U_s][\vec b/\vec y]
  \ty
  \Dia\bigl(C[\vec b/\vec y]\bigr).
\]
Thus, for the canonical source subterm, the generated elimination rule computes
to the original operational rewrite constructor.  Variables or other terms of
modal type may still be eliminated: their modal typing assumption is interpreted
as abstract evidence of the same interface, rather than as a syntactic claim that
the term literally has source shape \(U_s\).

\section{Interaction of \texorpdfstring{\(\Dia\)}{Diamond} with contextual modalities}
\label{sec:diamond-modal-interaction}

The possibility operator is a one-step operational effect, not a type-theoretic
counit.  The completion therefore does not contain a rule
\[
  \inferrule*{\Gamma\mid\Delta\vdash P\ty\Dia C}{\Gamma\mid\Delta\vdash P\ty C}.
\]
Such a rule would collapse may-reachability into actual membership and would
amount to a form of subject expansion.  The typed term calculus used in the
lambda-cube comparison is instead the constructor fragment: typed terms are
built from typed subterms by syntax-directed constructor rules.  The full modal
rules above remain available as a proof-theoretic layer for may-properties, but
unrestricted rewrite evidence is not a term former of the typed programming
calculus.

The interaction between \(\Dia\) and a contextual modality is \emph{admissible delayed}
elimination, and this admissible principle is available only when the relevant
site context has explicit closure support.  A judgment \(P\ty\Dia C\) says that \(P\) itself
has a one-step successor of type \(C\).  It does \emph{not} say that
\(K[P]\) can step inside every surrounding context \(K[-]\).  To move that step
through a particular site context one needs declared closure data
\(\kappa_s\Vdash K_s[-]\) in this sense.

Nested diamonds record the number of operational steps that have been postponed.
Write
\[
  \Dia^0 C = C,
  \qquad
  \Dia^{n+1}C = \Dia(\Dia^n C).
\]
For \(n=0\), modal elimination says that a term of contextual type can be
placed in its operational context and the result has type \(\Dia C\).  No
contextual-closure rule is needed in this case, because the step is the site
interaction itself.  For \(n>0\), a derivation of
\(P\ty\Dia^n M\) contains delayed rewrite evidence before the site interaction
can occur.  Each of those preliminary steps must be transported through
\(K_s[-]\) by an explicit closure support.

Concretely, for \(n>0\), delayed elimination is not added to the generated
presentation as a primitive rule.  It is the following admissible principle, available only when \(\kappa_s\Vdash K_s[-]\):
\begin{proposition}[Admissible delayed modal elimination]
\label{prop:admissible-delayed-modal-elim}
For every \(n>0\), if the site \(s\) has closure support
\(\kappa_s\Vdash K_s[-]\), then the following rule is admissible:
\begin{mathpar}
  \inferrule*[right={adm. delayed-modal-elim.}]
    {\kappa_s\Vdash K_s[-] \\
     \Gamma,\vec{x}:\vec X_s
       \mid \Delta,\vec e_x:\vec{x}\ty\vec{A}
       \vdash P\ty
       \Dia^n\bigl(\Mod{K_s}{\vec{x}\ty\vec{A}}{\vec{y}\ty\vec{B}}{C}\bigr) \\
     \Gamma,\vec{x}:\vec X_s
       \mid \Delta,\vec e_x:\vec{x}\ty\vec{A}
       \vdash \vec{b}\ty\vec{B}}
    {\Gamma,\vec{x}:\vec X_s
       \mid \Delta,\vec e_x:\vec{x}\ty\vec{A}
       \vdash K_s[P][\vec{b}/\vec{y}]\ty
       \Dia^{n+1}\bigl(C[\vec{b}/\vec{y}]\bigr)} .
\end{mathpar}
\end{proposition}
\begin{proof}[Proof sketch]
The proof is by induction on the displayed nested \(\Dia\)-evidence in the
premise.  For \(n=1\), the premise gives a one-step witness from \(P\) to some
\(P'\) with contextual type
\(\Mod{K_s}{\vec{x}\ty\vec{A}}{\vec{y}\ty\vec{B}}{C}\).  Closure support
transports that witness to a step
\(K_s[P][\vec b/\vec y]\rw K_s[P'][\vec b/\vec y]\), and ordinary modal
elimination applied to \(P'\) gives
\(K_s[P'][\vec b/\vec y]\ty\Dia(C[\vec b/\vec y])\).  A single use of
\(\Dia\)-introduction therefore gives the required
\(\Dia^2(C[\vec b/\vec y])\).  The successor step repeats the same transport
through \(K_s[-]\) before applying the induction hypothesis.  Thus the
unbounded family above is a metatheorem about nested possibility evidence, not
an infinite family of primitive rules in \(\mathcal H(A)\).
\end{proof}
The premise \(\kappa_s\Vdash K_s[-]\) is a side condition on the operational
presentation, not a consequence of the typing premise.  The successor case says:
if \(P\) may take \(n\) steps to a process with contextual type, then plugging
\(P\) into the site context may first propagate those delayed steps through
\(K_s[-]\), and then perform the site interaction itself.  The result is not a
term of type \(C\); it is a term of type \(\Dia^{n+1}C\).  Thus the admissible
principle composes may-information without ever erasing it.  In a proof-relevant
presentation, a derivation of the premise carries the staged rewrite evidence
being propagated.  In a proof-irrelevant predicate interpretation, the same
admissible principle uses the corresponding existential image and therefore
belongs to the regular structure used to interpret may-reachability, not to a
counit \(\Dia C\to C\).

For example, suppose the presentation contains left-parallel closure
\[
  r:P\rw P'
  \quad\longmapsto\quad
  \mathsf{par}_1(r):P\mid q \rw P'\mid q
\]
and the structural commutativity equation
\[
  P\mid q = q\mid P .
\]
Then closure support for the right-hole context \(q\mid\hole\) is obtained
by congruence transport:
\[
  q\mid P
  =
  P\mid q
  \rw
  P'\mid q
  =
  q\mid P' .
\]
Thus \(q\mid\hole\) need not be declared as a separate primitive closure
rule.  What commutativity does not do is create a new displayed site address
or a new profile quotient.  It only transports rewrite evidence that is
already available from the displayed closure rule.


This also explains how \(\Dia\) interacts with products in the lambda instance.
For \(n=0\), ordinary modal elimination gives
\[
  f\ty\dep(A,\lambda x.B),\quad a\ty A
  \quad\Longrightarrow\quad
  \app(f,a)\ty\Dia(B[a/x]).
\]
For \(n>0\), because the application context is closure-supported by the usual
head context rule, product elimination extends to
\begin{mathpar}
  \inferrule*[right={adm. delayed-\(\Pi\)-elim.}]
    {\mathsf{head}\Vdash\app(\hole,a) \\
     \Gamma\mid\Delta\vdash f\ty\Dia^n\bigl(\dep(A,\lambda x.B)\bigr) \\
     \Gamma\mid\Delta\vdash a\ty A}
    {\Gamma\mid\Delta\vdash \app(f,a)\ty\Dia^{n+1}\bigl(B[a/x]\bigr)} .
\end{mathpar}
In particular, from \(f\ty\Dia(\dep(A,\lambda x.B))\) and \(a\ty A\) one
gets
\[
  \app(f,a)\ty\Dia^2\bigl(B[a/x]\bigr),
\]
not \(\app(f,a)\ty B[a/x]\).  This is the desired behavior: the typing
judgment remembers that one step is needed to reach a function-like process and
a second step is needed to perform the application.  Without the explicit head
closure support, the admissible delayed application principle is not available.


\section{Conversion}

The completion does not identify types along operational rewrites.  In particular, it does not include a rule of the form
\[
\inferrule*{
  \Gamma\mid\Delta\vdash P \ty A
  \\
  \Gamma\mid\Delta\vdash r : A \rw B
  \\
  \Gamma\mid\Delta\vdash B \ty u
}{
  \Gamma\mid\Delta\vdash P \ty B
}
\]
for arbitrary rewrites \(r\).  Such a rule is appropriate only when the chosen rewrite relation is intended to be definitional equality at the type level, as in some presentations of beta-conversion for dependent lambda calculi.  It is not appropriate for a general operational theory, and especially not for concurrent calculi.

The only conversion principle present in the base construction is ordinary equational substitution.  If the presentation proves an equation \(A=B\), then
typing may be transported across that equation:
\[
\inferrule*{
  \Gamma\mid\Delta\vdash P \ty A
  \\
  \Gamma\mid\Delta\vdash A = B
}{
  \Gamma\mid\Delta\vdash P \ty B
}
\]
This is just the usual stability of judgments under the equations of the underlying presentation of the operational theory.

A particular instance may add a separate type-conversion relation \(\equiv\), together with a rule
\[
\inferrule*{
  \Gamma\mid\Delta\vdash P \ty A
  \\
  \Gamma\mid\Delta\vdash A \equiv B
  \\
  \Gamma\mid\Delta\vdash B \ty u
}{
  \Gamma\mid\Delta\vdash P \ty B
}
\]
but this is extra structure.

Likewise, subject reduction is not assumed by the generic construction.  A calculus may satisfy a theorem saying that \(P\ty A\) and \(P\rw Q\) imply \(Q\ty A\), or a suitable evolution of \(A\), but this is an additional property of a particular typed operational theory.  The modal introduction rule instead records the target typing explicitly.

\section{Church-style annotation of binding constructors}

The completion follows Barendregt's Church-style convention: constructors that bind variables are refined so that the type of the bound variable is part of the raw term.  The guiding principle is:
\begin{quote}
If a constructor discharges a typed variable in its generated typing rule, then the constructor receives an explicit type annotation for that variable.
\end{quote}

Suppose the input presentation has a binding constructor
\[
  c:Z\times\Proc^X\to\Proc.
\]
The completion contains a Church-style refinement
\[
  c^{\sharp}:Z\times X\times\Proc^X\to\Proc.
\]
The additional argument is the type annotation for the bound variable of carrier \(X\).  In this generic same-carrier formulation, an \(X\)-variable is typed by another element of \(X\); proof-relevantly, the generated typing display for that carrier lies over \(X\times X\).  Thus, in the \(\pi\)-calculus example, a name may also serve as a name-type or name-classifier.  A more refined presentation could replace the second \(X\) by a separate type-object \(\mathrm{Ty}_X\), but the present construction keeps the annotation carrier identical to the variable carrier. More generally, a constructor binding variables of carriers \(X_1,\ldots,X_m\) receives one type argument in each corresponding carrier.

There is an erasure map on raw syntax
\[
  \erase(c^{\sharp}(z,A,k))=c(z,k).
\]
The annotation is operationally irrelevant but syntactically relevant. Primitive rewrite constructors are lifted to their annotated versions by carrying the annotations through the source and target while erasing to the original rewrite.

For lambda abstraction, the raw constructor
\[
  \lam:\Proc^{\Proc}\to\Proc
\]
is refined to
\[
  \Lam:\Proc\times\Proc^{\Proc}\to\Proc.
\]
The beta rule is lifted to
\[
  A:\Proc,
  \quad
  t:\Proc^{\Proc},
  \quad
  y:\Proc
  \vdash
  \beta_A:\app(\Lam(A,t),y)\rw t(y).
\]

For asynchronous input, the raw constructor
\[
  \inp:\Nm\times\Proc^{\Nm}\to\Proc
\]
is refined to
\[
  \Inp:\Nm\times\Nm\times\Proc^{\Nm}\to\Proc,
\]
written \(n?_A.k\). The communication rule is lifted to
\[
  n,m,A:\Nm,
  \quad
  k:\Nm\to\Proc
  \vdash
  \comm_A:n!m\mid n?_A k\rw k(m).
\]

Non-binding constructors do not require annotations in the minimal Church-style lift. Application and output consume terms but do not bind them.


\section{Completing the untyped \texorpdfstring{$\lambda$}{λ}-calculus}

Start with the finite untyped lambda presentation with constructors
\[
  \app:\Proc\times\Proc\to\Proc,
  \qquad
  \lam:\Proc^{\Proc}\to\Proc,
\]
primitive beta rewrite
\[
  t:\Proc^{\Proc},
  \quad
  y:\Proc
  \vdash
  \beta:\app(\lam(t),y)\rw t(y),
\]
and the head contextual closure constructor
\[
  q:\Proc,
  \quad
  r:\Rew
  \vdash
  \mathsf{head}:\app(s(r),q)\rw\app(t(r),q).
\]
The head rule is not a product site; it is the explicit closure support for the
application context.  Use the product-site specification, which declares only
the function-position occurrence
\[
  U=\lam(t)
\]
in the beta source. The corresponding displayed context is
\[
  K[-]=\app(\hole,y).
\]
The variable \(y\) in this context
is a rely variable and there are no index variables.  The abstraction body \(t\)
is ambient: it occurs in the selected subterm \(\lam(t)\), not in the surrounding
application context.  In the familiar lambda rule this ambient parameter is the
schematic body of the abstraction, so it is displayed as the term typed in the
premise rather than as a new modal axis.  There is nothing to mod out by, so the
profile space for this site is \(\{\ast,\square\}^{2}\).  The argument occurrence
and the whole beta source are not declared product sites, so this completion
does not silently add argument-position or identity contextual modalities.

After Church refinement, the generated dependent product is notation for the corresponding contextual modality:
\[
  \dep(A, \lambda y.B)
  :=
  \Mod{\app(\hole,y)}{}{y\ty A}{B}.
\]
For a chosen profile \((u,v)\) in the selected subset \(\Sig\), the formation rule is
\begin{mathpar}
  \inferrule*[right={\(\Pi\)-form.}]
    {\Gamma\mid\Delta\vdash A\ty u \\
     \Gamma,y:\Proc\mid\Delta,e_y:y\ty A\vdash B\ty v}
    {\Gamma\mid\Delta\vdash \dep(A, \lambda y.B)\ty v}
\end{mathpar}

The abstraction rule is generated by the lifted beta rewrite and modal introduction:
\begin{mathpar}
  \inferrule*[right={\(\Pi\)-intro.}]
    {\Gamma\mid\Delta\vdash A\ty u \\
     \Gamma,y:\Proc\mid\Delta,e_y:y\ty A\vdash B\ty v \\
     \Gamma,y:\Proc\mid\Delta,e_y:y\ty A\vdash t\ty B}
    {\Gamma\mid\Delta\vdash \Lam(A, \lambda y.t)\ty\dep(A, \lambda y.B)}
\end{mathpar}
Indeed, under \(y:\Proc\mid e_y:y\ty A\), beta gives
\[
  \app(\Lam(A,\lambda y.t), y)\rw t,
\]
and the third premise gives \(t\ty B\). Modal introduction therefore gives
\[
  \Gamma\mid\Delta\vdash
  \Lam(A,\lambda y.t)
  \ty
  \Mod{\app(\hole,y)}{}{y\ty A}{B}.
\]

Application is modal elimination:
\begin{mathpar}
  \inferrule*[right={\(\Pi\)-elim.}]
    {\Gamma\mid\Delta\vdash f\ty\dep(A, \lambda x.B) \\
     \Gamma\mid\Delta\vdash a\ty A}
    {\Gamma\mid\Delta\vdash \app(f,a)\ty\Dia\bigl(B[a/x]\bigr)}
\end{mathpar}
or, using the built-in lambda syntax,
\begin{mathpar}
  \inferrule*[right={\(\Pi\)-elim.}]
    {\Gamma\mid\Delta\vdash f\ty\dep(A, k) \\
     \Gamma\mid\Delta\vdash a\ty A}
    {\Gamma\mid\Delta\vdash \app(f,a)\ty\Dia\bigl(k(a)\bigr)}
\end{mathpar}

So a function is a term that, when placed in the application context with an argument of type \(A\), can take one rewrite step to a result of type \(B\).
Because the lambda presentation includes the head contextual closure constructor,
Proposition~\ref{prop:admissible-delayed-modal-elim} also gives, for \(n>0\), the admissible delayed form
\begin{mathpar}
  \inferrule*[right={adm. delayed-\(\Pi\)-elim.}]
    {\mathsf{head}\Vdash\app(\hole,a) \\
     \Gamma\mid\Delta\vdash f\ty\Dia^n\bigl(\dep(A, \lambda x.B)\bigr) \\
     \Gamma\mid\Delta\vdash a\ty A}
    {\Gamma\mid\Delta\vdash \app(f,a)\ty\Dia^{n+1}\bigl(B[a/x]\bigr)}
\end{mathpar}
Thus a delayed function may still be applied by an admissible metarule, but each delayed operational step
remains visible as an additional \(\Dia\).

If the operational presentation is extended with an argument-position
closure rule
\[
  q:\Proc,\ r:\Rew
  \vdash
  \mathsf{arg}(q,r):
  \app(q,s(r))\rw \app(q,t(r)),
\]
then the argument-hole context \(\app(q,\hole)\) has closure support.  Hence,
in the nondependent case, there is an admissible strict delayed-argument
principle
\[
  \frac{
    f::A\to B
    \qquad
    g::\Dia^n A
  }{
    \app(f,g)::\Dia^{n+1}B
  }.
\]
This principle describes one possible reduction path: first reduce the
argument until it has type \(A\), then apply \(f\).  It is not a precise
normal-order cost rule.  Adding \(\mathsf{arg}\) enlarges the reduction
relation beyond pure normal-order reduction, and since \(\Dia\) is a
may-modality, the result records existence of such a path rather than the
least number of normal-order steps.

Moreover, no rule depending only on
\[
  f::A\to B
  \qquad
  g::\Dia A
\]
can determine the precise normal-order delay.  Functions with the same type
may discard, use, or duplicate their argument.  For example, if \(a::A\),
then
\[
  \app(\lambda x.a,g)\rw a
\]
gives a one-step possibility, while
\[
  \app(\lambda x.x,g)\rw g
\]
and \(g::\Dia A\) give a two-step possibility.  Thus precise delay accounting
requires additional usage or cost structure; it is not determined by the
ordinary function type alone.

\subsection{The lambda-cube geometry}

Since we require \((\ast,\ast)\) as a base profile, the remaining three optional profiles form the usual three-dimensional Lambda cube.

This is a recovery of the cube's sort geometry and constructor fragment, not an assertion that the completed modal calculus has exactly the same judgments as a pure type system.  The application rule is effectful: its result is \(\Dia(B[a/x])\), and delayed applications produce further nested diamonds. For comparison with Barendregt's systems we use the constructor-fragment translation of Section~\ref{sec:constructor-fragments-operational-properties}, which erases \(\Dia\).  We deliberately do not add a rule \(\Dia A\Rightarrow A\); that rule would type terms by subject expansion and would let arbitrary may-reachability masquerade as ordinary membership.

\section{Completing the asynchronous \texorpdfstring{$\pi$}{π}-calculus}

For the asynchronous communication rule
\[
  n,m:\Nm,
  \quad
  k:\Nm\to\Proc
  \vdash
  \comm:n!m\mid n?k\rw k(m),
\]
the displayed source is ordered with the output on the left and the input on the right.  This order is display data, even though the structural equations later identify it with the reversed parallel composite.  Use the receive-site specification, which declares the input occurrence
\[
  U=n?k.
\]
The corresponding context is
\[
  K[-]=n!m\mid\hole.
\]
This context is extracted from the displayed source tree.  For any plugged process \(P\), commutativity gives \(n!m\mid P=P\mid n!m\), but that equation does not create a second generated modality. The variable \(n\) appears both in \(U\) and in \(K\), so it is an index variable. The variable \(m\) appears in \(K\) but not in \(U\), so it is a rely variable. The continuation \(k\) appears in the subterm and in the target \(k(m)\), but not in the surrounding context, so it is an ambient variable:
\[
  \vec{x}=(n),
  \qquad
  \vec{y}=(m),
  \qquad
  \vec{z}=(k),
  \qquad
  \Theta_{\mathsf{recv}}=(k:\Nm\to\Proc).
\]
The declaration for \(k\) is displayed as a structural higher-order declaration; if the continuation variables themselves carry a typed discipline, that additional information belongs in the same ambient telescope.

After Church refinement, input is written
\[
    \Inp(n, B, k),
\]
or with syntactic sugar as
\[
  n?_B k,
\]
where \(B\) classifies the received message.  For a chosen profile \((\alpha,\beta;\gamma)\), assigning sorts to the index type \(A\) of the channel \(n\), the rely type \(B\) of the message \(m\), and the target type \(C\), the generated input rule uses the ambient, index, and rely carrier context
\[
  \Gamma,\Theta_{\mathsf{recv}},n:\Nm,m:\Nm
\]
together with explicit relation-context evidence assumptions
\[
  \Delta,e_n:n\ty A,e_m:m\ty B .
\]
The rely variable \(m\) and its evidence assumption \(e_m:m\ty B\) are
discharged by the modal type former, while the index variable \(n\) and
its evidence assumption \(e_n:n\ty A\) remain in the external context.
\begin{mathpar}
  \inferrule*[right={input.}]
    {
      \Gamma,\Theta_{\mathsf{recv}}\mid\Delta
      \vdash A\ty\alpha
      \\
      \Gamma,\Theta_{\mathsf{recv}},n:\Nm
      \mid
      \Delta,e_n:n\ty A
      \vdash B\ty\beta
      \\
      \Gamma,\Theta_{\mathsf{recv}},n:\Nm,m:\Nm
      \mid
      \Delta,e_n:n\ty A,e_m:m\ty B
      \vdash C\ty\gamma
      \\
      \Gamma,\Theta_{\mathsf{recv}},n:\Nm,m:\Nm
      \mid
      \Delta,e_n:n\ty A,e_m:m\ty B
      \vdash k(m)\ty C
    }
    {
      \Gamma,\Theta_{\mathsf{recv}},n:\Nm
      \mid
      \Delta,e_n:n\ty A
      \vdash
      n?_B k
      \ty
      \Mod{n!m\mid\hole}{n\ty A}{m\ty B}{C}
    }
\end{mathpar}
The elimination rule keeps the same ambient continuation parameter:
\begin{mathpar}
  \inferrule*[right={input-elim.}]
    {
      \Gamma,\Theta_{\mathsf{recv}},n:\Nm \;|\;
      \Delta,e_n:n\ty A
      \vdash
      P\ty\Mod{n!m\mid\hole}{n\ty A}{m\ty B}{C}
      \\
      \Gamma,\Theta_{\mathsf{recv}},n:\Nm \;|\;
      \Delta,e_n:n\ty A
      \vdash a\ty B
    }
    {
      \Gamma,\Theta_{\mathsf{recv}},n:\Nm \;|\;
      \Delta,e_n:n\ty A
      \vdash
      n!a\mid P\ty\Dia\bigl(C[a/m]\bigr)
    }
\end{mathpar}
Thus an input process has a type describing the operational guarantee that, when paired with a suitable output, it has a one-step possibility of reaching a continuation of type \(C\).  

The output occurrence \(n!m\) may be declared as a separate site to generate the dual output-context modality, but it is not part of the receive-site hypercube below; it would produce its own hypercube.

\subsection{The receive hypercube.}

For the receive site \(s_{\mathsf{recv}}\), no consistency relation identifies
or rules out any profiles, so the carrier-sensitive profile set is exactly
\[
  \{\ast^{\Nm},\square^{\Nm}\}_{n}
  \times
  \{\ast^{\Nm},\square^{\Nm}\}_{m}
  \times
  \{\ast^{\Proc},\square^{\Proc}\}_{C}.
\]
When the carriers are clear, we abbreviate this as \(\{\ast,\square\}^{3}\). Since we require the base profile
\((\ast^{\Nm},\ast^{\Nm};\ast^{\Proc})\), the remaining seven profiles are optional, so there are \(2^7=128\) based signature choices.  A profile \((\alpha,\beta;\gamma)\in \Sigma \subseteq \{\ast,\square\}^{3}\), with carrier superscripts suppressed when unambiguous, should be read as follows: \(\alpha\) is the sort of the channel type \(A\) for the indexed channel \(n\), \(\beta\) is the sort of the message type \(B\) for the received name \(m\), and \(\gamma\) is the sort of the continuation target type \(C\).  The base profile \((\ast^{\Nm},\ast^{\Nm};\ast^{\Proc})\) gives first-order receive: ordinary channels receive ordinary messages and continue at an ordinary process type.  The profile \((\ast^{\Nm},\ast^{\Nm};\square^{\Proc})\) is the receive analogue of dependent types: a type-level protocol family may depend on ordinary channel and message values. Profiles with \(\gamma=\ast^{\Proc}\) and at least one of \(\alpha,\beta\) equal to \(\square^{\Nm}\) are polymorphic receive principles: ordinary continuations may be indexed by type-level channels, by type-level payloads, or by both.  Profiles with \(\gamma=\square^{\Proc}\) and at least one of \(\alpha,\beta\) equal to \(\square^{\Nm}\) are receive type-operator principles: protocol families may themselves be parameterized by type-level channels, type-level payloads, or both.  Finally, because the message type \(B\) is formed in the context\(n\ty A\), the pair \((\alpha,\beta)\) also controls the dependency of the message classifier on the channel classifier: \((\ast^{\Nm},\ast^{\Nm})\) is the ordinary first-order case, \((\ast^{\Nm},\square^{\Nm})\) allows channel-dependent message classifiers, and profiles with \(\alpha=\square^{\Nm}\) allow the channel itself to be indexed by type-level name data.

\section{Freeness and functoriality}
\label{sec:functoriality-adjunction}

The construction is not merely a recipe.  It is free: once the displayed sites,
profile signatures, and closure-support annotations have been chosen, no hidden
typing structure is imposed.  To interpret the completed system in another typed
operational theory is exactly to interpret the underlying site-and-signature
data, together with the chosen generated structure.

The precise categories are defined in Appendix~\ref{app:functoriality-adjunction}.
Here \(\mathbf{PSig}\) is the category of finite site-and-signature
presentations and strict morphisms preserving displayed source addresses,
matching data, signatures, and closure support.  The target category
\(\mathbf{PDTOT}\) consists of typed operational presentations equipped with a
specified interpretation of the generated hypercube structure.  The forgetful
functor
\[
  U:\mathbf{PDTOT}\to\mathbf{PSig}
\]
forgets that interpretation.

\begin{theorem}[Free hypercube completion]
\label{thm:free-hyp}
The structured completion functor
\[
  \Hyp:\mathbf{PSig}\to\mathbf{PDTOT}
\]
is left adjoint to \(U\).  Equivalently, for every
\(A\in\mathbf{PSig}\) and every structured typed operational theory
\(X\in\mathbf{PDTOT}\), there is a natural bijection
\[
  \mathrm{Hom}_{\mathbf{PDTOT}}(\Hyp(A),X)
  \cong
  \mathrm{Hom}_{\mathbf{PSig}}(A,U(X)).
\]
\end{theorem}

\begin{proof}[Proof sketch]
A morphism out of \(\Hyp(A)\) is forced by its action on the underlying
site-and-signature presentation \(A\).  The sorts, typing evidence displays,
modal constructors, source-subterm rules, and computation equations of
\(\Hyp(A)\) are freely generated from that data.  Closure-support annotations are part of the underlying site data; they justify admissible delayed elimination but do not add primitive delayed rules.  Conversely, any strict
morphism \(A\to U(X)\) extends uniquely by applying the structure map of \(X\)
to each generated symbol and rule.  Thus the displayed hom-set map sends a
structured morphism to its underlying \(\mathbf{PSig}\)-morphism, and the inverse
extends a \(\mathbf{PSig}\)-morphism by freeness.  The unit is the identity on
\(A\), and the counit is the specified structure map of \(X\).  The triangle
identities reduce to identity laws.  Appendix~\ref{app:functoriality-adjunction}
contains the formal category definitions, functoriality proof, unit, counit, and
triangle verification.
\end{proof}




\section{Proof-relevant evidence semantics and soundness}
\label{sec:soundness}

The generated rules have a proof-relevant semantic reading.  A typing judgment is
not interpreted first as a truth-valued predicate; it is interpreted as an
object of evidence.  For each relevant carrier \(X\), a structured model carries
a display
\[
  p_X:\Ev_X\longrightarrow X\times X,
  \qquad
  p_X=\langle\tm_X,\tp_X\rangle,
\]
not assumed to be monic.  An element \(e\in\Ev_X\) is evidence that
\(\tm_X(e)\) has classifier \(\tp_X(e)\).  Thus a derivation of
\(\Gamma\mid\Delta\vdash P\ty A\) is interpreted by a morphism
\[
  e_P:\llbracket\Gamma\mid\Delta\rrbracket\longrightarrow \Ev_X
\]
whose projection to \(X\times X\) is
\(\langle\llbracket P\rrbracket_{\Gamma\mid\Delta},
\llbracket A\rrbracket_{\Gamma\mid\Delta}\rangle\).

Typed context extension is still interpreted by pullback.  If
\(A:X\) is a type code in context \(\Gamma\mid\Delta\), then
\[
  \llbracket\Gamma,x:X\mid\Delta,e_x:x\ty^X A\rrbracket
  =
  (\llbracket\Gamma\mid\Delta\rrbracket\times X)
  \times_{X\times X}
  \Ev_X,
\]
where the map \(\llbracket\Gamma\mid\Delta\rrbracket\times X\to X\times X\)
is \(\langle\pi_2,\llbracket A\rrbracket\pi_1\rangle\).  A point of this
pullback is a carrier element together with named evidence witnessing that it
has the displayed type.  Weakening and substitution are therefore ordinary
pullback operations on these evidence displays.

With this convention, the modal rules retain their operational witnesses.
Evidence for \(P\ty\Dia C\) is one-step may-evidence: it supplies a reduct
\(Q\), a rewrite witness \(r:P\rw Q\), and evidence that \(Q\ty C\).  The
\(\Dia\)-introduction rule constructs this tuple, and the \(\Dia\)-eliminator
is its dependent eliminator.  Similarly, evidence for a contextual modality is
an interface for the declared site context: when supplied with rely arguments
and their evidence, modal elimination returns \(\Dia\)-evidence for the plugged
process.  Source-subterm introduction is sound because the primitive rewrite
constructor itself supplies the required one-step witness.

\begin{theorem}[Soundness]
\label{thm:soundness}
In every structured proof-relevant model of the generated typed operational
presentation, every derivable typing judgment has semantic evidence.
\end{theorem}

\begin{proof}[Proof sketch]
The proof is by induction on derivations.  Structural rules are interpreted by
finite-limit operations on evidence contexts.  Rules that discharge typed
variables, such as modal introduction and Church-style abstraction, use the
specified binding/exponential structure of the operational presentation, or
equivalently the named rule morphisms supplied by a structured model.  The
\(\Dia\)-introduction rule pairs a rewrite witness with evidence for the typed
target.  The \(\Dia\)-eliminator destructs this proof-relevant evidence.  Modal
elimination applies the evidence interface carried by a contextual modality to
supplied rely evidence.  The admissible delayed-elimination principle is sound
only for sites with explicit closure support; the closure constructor transports
the staged rewrite evidence through the displayed context before the site
interaction occurs.  Appendix~\ref{app:soundness} gives the rule-by-rule
verification.
\end{proof}

\begin{proposition}[Proof-irrelevant may semantics]
\label{prop:regular-may}
The proof-relevant soundness theorem does not require typing displays to be
subobjects.  If one forgets evidence and wants to interpret \(\Dia C\) and
contextual modalities as proof-irrelevant may-predicates, then extra structure
is required: regular images for one-step possibility, and dependent universal
operations for contextual modalities that discharge rely variables.
\end{proposition}

\begin{proof}[Proof sketch]
Proof-relevant judgments retain rewrite and typing evidence.  Proof-irrelevant
may semantics forgets that evidence and asks only whether some evidence exists;
this existential collapse is an image operation.  For contextual modalities, the
rely variables are abstracted out of the interface, so the proof-irrelevant
reading also needs the corresponding dependent operation over those variables.
Appendix~\ref{app:soundness} spells out this optional evidence-forgetting
semantics.
\end{proof}


\section{Constructor fragments and operational properties}
\label{sec:constructor-fragments-operational-properties}

The completion generates modal typing rules from operational sites.  These
rules have a direct may-reading, but they should not all be confused with
syntax-directed term formation rules.  In the typed programming language, the
term calculus we use is the \emph{constructor fragment}; the full completed
system is a specification layer for may-properties of operational syntax.

By the constructor fragment we mean the subsystem in which typed terms are
built only by the generated typed constructors and their syntax-directed
rules.  For the lambda-calculus completion, this includes the sort and
structural rules, product formation, the \(\Lam\)-introduction rule, the
\(\app\)-elimination rule, and the admissible delayed application principles justified by the explicitly supplied
head closure support.  These admissible principles add no primitive typing rules to the generated presentation.  The fragment does not include unrestricted
\(\Dia\)-introduction or modal-introduction applied to arbitrary raw terms.
Thus every immediate object-language subterm of a typed process is itself typed
by a premise of the rule that constructs it, even when the resulting type
contains one or more \(\Dia\)'s.

\subsection{Termination in the \texorpdfstring{$\lambda$}{λ}-calculus}

If a term is well-typed in one of the calculi of the lambda
cube, it necessarily terminates.  Consider the completion of the
untyped lambda calculus in which we keep only the dependent-
product-like site
\[
  K[-]=\app(\hole,x).
\]
For a based selected set of product profiles
\(\Sig\subseteq\{\ast,\square\}^2\), with \((\ast,\ast)\in\Sig\), and for each
\((u,v)\in\Sig\), the constructor fragment has the following primitive rules,
together with the displayed admissible delayed principle:
\begin{mathpar}
  \inferrule*[right={\(\Pi\)-form}]
    {\Gamma\mid\Delta\vdash A\ty u \\
     \Gamma,x:\Proc\mid\Delta,e_x:x\ty A\vdash B\ty v}
    {\Gamma\mid\Delta\vdash \dep(A,\lambda x.B)\ty v}

  \inferrule*[right={\(\Pi\)-intro}]
    {\Gamma\mid\Delta\vdash A\ty u \\
     \Gamma,x:\Proc\mid\Delta,e_x:x\ty A\vdash B\ty v \\
     \Gamma,x:\Proc\mid\Delta,e_x:x\ty A\vdash t\ty B}
    {\Gamma\mid\Delta\vdash \Lam(A,\lambda x.t)\ty\dep(A,\lambda x.B)}

  \inferrule*[right={\(\Pi\)-elim}]
    {\Gamma\mid\Delta\vdash f\ty\dep(A,\lambda x.B) \\
     \Gamma\mid\Delta\vdash a\ty A}
    {\Gamma\mid\Delta\vdash \app(f,a)\ty\Dia\bigl(B[a/x]\bigr)}

  \inferrule*[right={adm. delayed-\(\Pi\)-elim, \(n>0\)}]
    {\mathsf{head}\Vdash\app(\hole,a) \\
     \Gamma\mid\Delta\vdash f\ty\Dia^n\bigl(\dep(A,\lambda x.B)\bigr) \\
     \Gamma\mid\Delta\vdash a\ty A}
    {\Gamma\mid\Delta\vdash \app(f,a)\ty\Dia^{n+1}\bigl(B[a/x]\bigr)}
\end{mathpar}
The \(\Pi\)-elimination rule is the generated modal elimination rule.  The cases
\(n>0\) are admissible delayed eliminations justified by the explicitly supplied head
closure support, not primitive members of the generated presentation.  The result type is therefore \(\Dia^{n+1}(B[a/x])\), not
\(B[a/x]\).  Nevertheless, in the constructor fragment this difference does
not introduce nonterminating untyped subterms.

Define a translation \((-)^\circ\) from types and terms in this fragment to the corresponding calculus in the lambda cube by erasing every possibility modality and by sending the generated product, application, and abstraction constructors to their counterparts:
\[
  (\Dia^n C)^\circ = C^\circ,
\]
\[
  \bigl(\dep(A,\lambda x.B)\bigr)^\circ
  =
  \Pi_{x:A^\circ}B^\circ,
\]
\[
  \bigl(\app(M,N)\bigr)^\circ
  =
  (M^\circ \; N^\circ),
\]
\[
  \bigl(\Lam(A,\lambda x.t)\bigr)^\circ
  =
  \lambda x:A^\circ.\,t^\circ.
\]

\begin{proposition}[Termination in the lambda constructor fragment]
Let \(P\) be a process term typable in the constructor fragment of the product-site completion of the untyped lambda calculus.  Then \(P\) is strongly normalizing for beta-reduction.
\end{proposition}

\begin{proof}[Proof sketch]
The proof is by translation to the corresponding Barendregt system.  By
induction on the typing derivation, if
\[
  \Gamma\mid\Delta\vdash P\ty A
\]
is derivable in the constructor fragment, then
\[
  \Gamma^\circ\vdash P^\circ : A^\circ
\]
is derivable in the corresponding pure type system.  The only nontrivial case
is the admissible delayed application principle.  That principle concludes
\[
  \Gamma\mid\Delta\vdash \app(f,a)\ty\Dia^{n+1}\bigl(B[a/x]\bigr)
\]
from premises \(f\ty\Dia^n(\dep(A,\lambda x.B))\) and \(a\ty A\).
After erasing all diamonds, this becomes exactly the usual Barendregt
application conclusion
\[
  \Gamma^\circ\vdash
  (f^\circ \; a^\circ)
  :
  B^\circ[a^\circ/x].
\]
By the standard strong-normalization theorem for the lambda cube, in particular for the calculus of constructions \cite{Barendregt1991,Barendregt1992Types}, every
constructor-fragment subsystem obtained by choosing fewer product profiles is strongly
normalizing.  The translation preserves the computational beta-redexes, so
termination of \(P^\circ\) implies termination of \(P\).  If the operational
presentation uses a sub-strategy such as normal-order beta reduction, this
follows a fortiori from strong normalization for full beta reduction.
\end{proof}

The restriction to the constructor fragment is essential.  The full modal
completion also has rules that type a process from the existence of a rewrite
witness.  Those rules express may-behavior, but they are not syntax-directed
term formation rules and can type processes whose subterms were not built
from typed components.  The termination statement above is therefore a
statement about the lambda-cube-like fragment, not about the entire completed
modal theory.

\subsection{The analogous reading for the \texorpdfstring{\(\pi\)}{π}-calculus}

For a concurrent calculus, the analogous theorem is not necessarily a termination theorem. A process may be intended to run forever, wait for messages, participate in a race, or interact with an environment.  The generated modal types instead describe one-step reachability properties.

Suppose \(N:\Proc\) is a type intended to classify the nil process, and
suppose the model or typing presentation contains
\begin{mathpar}
  \inferrule*[right={nil.}]
    { }
    {\Gamma\mid\Delta\vdash {\rm nil}\ty N}
\end{mathpar}
If \(N\) is interpreted as the property of being structurally congruent to
\({\rm nil}\), then \(P\ty\Dia N\) says that \(P\) may take one step to the
nil process.  More generally, the rule
\begin{mathpar}
  \inferrule*[right={\(\Dia\)-intro.}]
    {\Gamma\mid\Delta\vdash r:P\rw Q \\
     \Gamma\mid\Delta\vdash Q\ty N}
    {\Gamma\mid\Delta\vdash P\ty\Dia N}
\end{mathpar}
says that \(P\) has an explicit one-step witness to some \(N\)-typed target.

The receive site gives the contextual version of the same idea.  From the
communication rule
\[
  n!m\mid n?k \rw k(m),
\]
the generated receive typing rule includes the instance, written with the
ambient continuation telescope \(\Theta_{\mathsf{recv}}=(k:\Nm\to\Proc)\):
\begin{mathpar}
  \inferrule*[right={recv-nil.}]
    {
      \Gamma,\Theta_{\mathsf{recv}},n:\Nm,m:\Nm
      \mid
      \Delta,e_n:n\ty A,e_m:m\ty B
      \vdash k(m)\ty N
    }
    {
      \Gamma,\Theta_{\mathsf{recv}},n:\Nm
      \mid
      \Delta,e_n:n\ty A
      \vdash
      n?_B k
      \ty
      \Mod{n!m\mid\hole}{n\ty A}{m\ty B}{N}
    }
\end{mathpar}
Thus a process of type
\[
  \Mod{n!m\mid\hole}{n\ty A}{m\ty B}{N}
\]
is one which, when placed in parallel with an output \(n!m\) supplying a
message \(m\) with evidence \(e_m:m\ty B\) along an indexed channel \(n\) with evidence \(e_n:n\ty A\), may take one
communication step to a process of type \(N\).  If \(N\) is the nil property,
this says that the process may communicate once and reach nil---but it may do something else!

The elimination rule makes this operational reading explicit:
\begin{mathpar}
  \inferrule*[right={recv-elim}]
    {
      \Gamma,\Theta_{\mathsf{recv}},n:\Nm
      \mid
      \Delta,e_n:n\ty A
      \vdash
      P\ty\Mod{n!m\mid\hole}{n\ty A}{m\ty B}{N}
      \\
      \Gamma,\Theta_{\mathsf{recv}},n:\Nm
      \mid
      \Delta,e_n:n\ty A
      \vdash a\ty B
    }
    {
      \Gamma,\Theta_{\mathsf{recv}},n:\Nm
      \mid
      \Delta,e_n:n\ty A
      \vdash
      n!a\mid P\ty\Dia\bigl(N[a/m]\bigr)
    }
\end{mathpar}
When \(N\) does not depend on \(m\), the conclusion is simply
\[
  n!a\mid P\ty\Dia N.
\]

Nested modalities express staged may-reachability.  For example, a type of
the form
\[
  \Mod{K_2}{x_2\ty A_2}{y_2\ty B_2}
    {\Mod{K_1}{x_1\ty A_1}{y_1\ty B_1}{N}}
\]
says that, after the outer index data are fixed and the outer context \(K_2\) and its rely data are supplied,
the process may step to a new process satisfying the inner modal type; after
the inner index data are fixed and the inner context \(K_1\) and its rely data are supplied, that continuation
may step to a process of type \(N\).  If \(N\) denotes nil, this is a
two-stage may-reach-nil property.

\subsection{Must properties are model-theoretic}

The modal types generated by the completion have an existential, proof-relevant
shape.  A derivation of
\[
  P\ty\Dia C
\]
contains, or is generated from, a particular rewrite witness
\[
  r:P\rw Q
\]
and a typing derivation
\[
  Q\ty C.
\]
Similarly, a derivation of
\[
  P\ty\Mod{K}{\vec{x}\ty\vec{A}}{\vec{y}\ty\vec{B}}{C}
\]
records that, after the index variables are fixed and the rely variables are supplied, the plugged process
\(K[P]\) has a displayed one-step path to a \(C\)-typed target.  Any ambient
variables of the selected source subterm remain visible in the surrounding
context throughout this derivation; they are not hidden in the modal interface.
These are may-properties.

A must-property has a different quantifier shape.  One-step must preservation
would say something like
\[
  \forall Q.\;
  K[P]\rw Q
  \Longrightarrow
  Q\in\llbracket C\rrbracket.
\]
Such statements quantify over all rewrites in a model \(\mathcal M\).  They are therefore not determined by the syntactic modal typing rules alone.

This distinction matters especially for concurrent calculi.  An arbitrary model of a presented operational theory may contain additional elements, additional rewrite witnesses, additional nondeterminism, or additional identifications beyond the freely generated syntax.  A may-typing derivation is stable under such enlargement: the displayed witness still exists.  A must-claim is not stable in the same way, because adding a new possible rewrite can invalidate a universal statement about all rewrites.

The free model is the special case in which the only processes, equations, and rewrite witnesses are those generated by the presentation, modulo the specified equations.  In the free model it can make sense to prove must-style metatheorems by induction on generated syntax or by operational analysis.  Strong normalization for the lambda-cube constructor fragment is one such theorem.  For a \(\pi\)-calculus fragment, analogous free-model theorems might say that every generated run respecting a closed protocol eventually reaches a process of type \(N\), or that every generated communication step preserves some protocol invariant.  But these are theorems about a chosen model or class of models, not inference rules supplied by the hypercube completion itself.

\section{Relation to other presentations}

\paragraph{Displayed evidence and subobject fibrations.}
The main semantics uses displayed evidence objects rather than mere truth-valued
subobjects: a typing judgment is interpreted by a section of a display over the
corresponding term-and-classifier pair.  Finite limits provide the pullbacks
needed for typed context extension, weakening, substitution, and the premise
objects of generated rules.  If one forgets evidence, the resulting
proof-irrelevant predicates live in a subobject fibration \cite{Jacobs1999,LambekScott1986}; existential images, and for contextual modalities with discharged rely variables the corresponding dependent universal operations, are then needed for the optional proof-irrelevant interpretation of may-reachability.

\paragraph{Generalized algebraic theories.}
The generated rules are algebraic in the contextual sense: typed contexts are finite-limit objects, eliminators are morphisms out of those objects, and computation rules are equations.  This is closely aligned with generalized algebraic theories and contextual categories \cite{Cartmell1986}.

\paragraph{Binding.}
The use of specified exponentials for continuations follows the abstract-syntax tradition in which binding structure is supplied as presentation data \cite{FiorePlotkinTuri1999}.  Because the exponential \(\Proc^\Nm\) is specified in the asynchronous \(\pi\)-calculus presentation, its models are required to preserve the corresponding evaluation and substitution structure.

\paragraph{Reactive systems.}
Representing rule instances and contexts explicitly is compatible with the general lesson of reactive systems: operational behavior is organized by structured contexts and evidence, rather than by an unstructured binary relation alone \cite{LeiferMilner2000}.

\section{Conclusion}

The guiding observation of this paper is that dependent products are not isolated artifacts of lambda syntax.  They are generated by a particular operational site: the function position of the beta redex.  Once this is made explicit, the lambda cube becomes one instance of a more general construction.  A displayed rewrite source supplies possible sites; a site supplies a one-hole context; and a choice of quotient sort profiles determines which modal dependent type formers are present.

The construction is deliberately presentation-sensitive.  Displayed syntax names the sites, structural equations govern equality in the operational theory, and profile-matching data computes the quotient sort geometry.  These layers are kept separate so that, for example, commutativity of parallel composition does not silently move holes or merge input and output sites.  Likewise, the possibility modality records one-step may-evidence; it does not include a counit \(\Dia A\Rightarrow A\), and delayed use of a modal type requires explicit contextual-closure support.

For the lambda-calculus, the application site recovers Barendregt's product geometry in the constructor fragment.  For the asynchronous \(\pi\)-calculus, the communication rule gives different local cubes, with name and process carriers contributing their own sort levels.  Thus the same recipe generates a family of typed operational systems from the redex structure of the language, rather than from lambda application alone.

The formal results support this recipe.  The generated-presentation assignment is functorial, the structured completion is free in the sense of the adjunction proved in Appendix~\ref{app:functoriality-adjunction}, and the generated typing rules are sound in the proof-relevant evidence semantics proved in Appendix~\ref{app:soundness}.  The main text has kept the examples and intuitions in front; the appendices contain the finite profile algorithm, the generated rule specification, and the full categorical and semantic verifications.

\clearpage
\appendix

\section{Formal site-marked presentations and profile quotients}
\label{app:site-profile-details}
This appendix gives the formal version of the site and profile construction used in Section~\ref{sec:construction}.
As noted above, a site is a selected occurrence in the displayed source of a rewrite rule with carrier $\Proc$ that gets replaced with a hole to form a term context.  When the selected subterm is placed into the context, a reduction may occur.  The selection is deliberate: not every displayed $\Proc$-subterm of a rewrite source should automatically become a modal type former.  Sites are named before quotienting by structural equations; the completion is therefore presentation-sensitive in the same way that an operational semantics written with evaluation contexts is presentation-sensitive.

\subsection{The fixed sorts}

For each relevant carrier object \(X\), the generated presentation adjoins constants
\[
  \ast^X:X,
  \qquad
  \square^X:X,
\]
for types and kinds, respectively, and a proof-relevant typing-evidence
display
\[
  p_X:\Ev_X\longrightarrow X\times X,
  \qquad
  p_X=\langle\tm_X,\tp_X\rangle .
\]
A point of \(\Ev_X\) is evidence that the process or type code
\(\tm_X(e)\) has classifier \(\tp_X(e)\).  In contexts we write this evidence
with a name, as \(e:a\ty^X b\).  When the evidence name is immaterial in the
conclusion of a rule, we may still write \(a\ty^X b\).  The sort axiom is
parametric in \(X\):
\begin{mathpar}
  \inferrule*[right={sort.}]
    { }
    {\Gamma\mid\Delta\vdash \ast^X\ty^X\square^X}
\end{mathpar}
When the superscript is clear from the context, it is usually suppressed.


\subsection{Displayed syntax, structural equations, and profile constraints}

The construction is deliberately presentation-sensitive.  It distinguishes the raw displayed syntax tree of a primitive rewrite source from the structural congruence generated by equations.  A primitive rewrite constructor
\[
  \Gamma_i,\Delta_i\vdash \rho_i:L_i\rw R_i
\]
therefore supplies a named displayed redex representative \(L_i\rw R_i\).  Sites are occurrence addresses in the raw tree \(L_i\), before quotienting by any equations of the presentation.

Displayed equations have two different possible roles, and these roles are kept separate.  First, they are ordinary equations of the operational theory: they generate structural congruence and support equational substitution in judgments.  Second, an equation schema may carry explicit type-occurrence matching data \(\Match_s(E)\), used only by the profile-consistency algorithm.  The first role does not imply the second.  In particular, an equation such as commutativity or associativity of a constructor does not automatically identify sites, replace a context by all congruent contexts, or identify the sort slots attached to moved occurrences.

Thus the following data are independent inputs to the completion:
\begin{enumerate}[label=\textup{(\alph*)}]
  \item the displayed source representative \(L_i\) of each primitive rewrite constructor;
  \item the finite site specification \(\mathfrak S_i\), whose elements are occurrence addresses in that displayed source;
  \item the displayed equations of the presentation, used as ordinary equations;
  \item the finite matching relations \(\Match_s(E)\), used to generate profile-slot identifications.
\end{enumerate}
Changing the displayed representative of a redex, adding an equivalent representative as a separate primitive rewrite, or adding explicit equation matchings can change the generated modal operators.  Merely adding a structural equation changes the equational theory but not the set of sites extracted from a fixed source display.


\subsection{Process sites}

Let
\[
  \Gamma_i,\Delta_i\vdash \rho_i:L_i\rw R_i
\]
be a primitive rewrite constructor in a finite presentation \(\mathcal{P}\). Since
\(L_i\) is a finite raw syntax tree, it has finitely many displayed subterm
occurrences, even if \(L_i\) is structurally congruent to many other terms.  Write \(\epsilon_i\) for the root occurrence of \(L_i\), and write
\[
  \Occ_{\Proc}(L_i)
\]
for the finite set of occurrences whose displayed subterm has carrier \(\Proc\).
Let
\[
  \Occ^{\mathrm{ch}}_{\Proc}(L_i)
  \subseteq
  \Occ_{\Proc}(L_i)
\]
be the subset of proper occurrences \(o\ne\epsilon_i\) whose displayed subterm
is headed by an object-language constructor of \(\mathcal{P}\).  Variable
occurrences and occurrences headed by the structural source and target maps
\(s,t:\Rew\to\Proc\) are not in \(\Occ^{\mathrm{ch}}_{\Proc}(L_i)\).

The construction does not use all of \(\Occ_{\Proc}(L_i)\) by default.  Instead,
the operational presentation carries a finite \emph{site specification}
\(\mathfrak{S}\), a family
\[
  \mathfrak{S}_i\subseteq \Occ_{\Proc}(L_i)
\]
choosing the source positions that are to generate contextual modal types.  If a
bare presentation is used without an explicit site specification, the default
site policy in this paper is the finite family
\[
  \mathfrak{S}^{\mathrm{def}}_i
  =
  \Occ^{\mathrm{ch}}_{\Proc}(L_i).
\]
Thus the default chooses proper constructor-headed process occurrences and
excludes variable occurrences, structural \(s,t\)-occurrences, and the whole
rewrite source.  This default is only a convention for examples; a language
designer may declare a smaller or larger finite set.  In particular, variable
occurrences and the whole source occurrence generate modalities only when they
are explicitly declared.  The separate possibility type \(\Dia\) below covers
the identity-hole case, so the whole source is not included by default.

For example, in the $\pi$-calculus, the comm rule says
\[
    n,m: \Nm, k:\Nm \to \Proc \vdash {\rm comm}: n!m \;|\; n?k \rw k(m).
\]
The proper constructor-headed source occurrences with carrier \(\Proc\) are
\(n!m\) and \(n?k\).  The default site policy declares both.  The input occurrence \(n?k\), with displayed context \(K[-]=n!m\mid[-]\), gives the receive modality
\[
    \langle n!m \;|\; [-]\rangle^{n\ty A}_{m\ty B}C.
\]
The output occurrence has the dual displayed context \(K[-]=[-]\mid n?k\), but its interface is not the simple expression \(k\ty A\): it involves the higher-order continuation parameter \(k:\Nm\to\Proc\) and a different ambient/rely/index split from the receive site.  We therefore leave the dual modality in schematic site notation unless an explicit type discipline for continuations has been chosen.  These are declarations about the two children of the displayed source tree \(n!m\;|\;n?k\).  The commutativity equation for \(\mid\) does not collapse them into one site and does not add the reversed-site contexts \(n?k\;|\;[-]\) or \([-]\;|\;n!m\).  Declaring only the input occurrence gives the receive fragment used later; declaring neither adds no communication-specific contextual modality.

\begin{definition}[Site specification and process sites]
A site-marked operational presentation is a pair \((\mathcal{P},\mathfrak{S})\)
where \(\mathcal{P}\) is a finite operational presentation and \(\mathfrak{S}\)
is a finite site specification as above.  Its process-site set is
\[
  \Sites(\mathcal{P},\mathfrak{S})
  =
  \{(i,o)\mid o\in\mathfrak{S}_i\}.
\]
When the site specification is fixed, we write \(\Sites(\mathcal{P})\) for
\(\Sites(\mathcal{P},\mathfrak{S})\).  For \(s=(i,o)\in\Sites(\mathcal{P})\),
write
\[
  U_s:\Proc
\]
for the selected displayed subterm occurrence. There is a displayed one-hole
process context \(K_s[-]\) such that
\[
  L_i = K_s[U_s].
\]
This is equality of displayed syntax trees, not equality modulo the equations of the operational theory.  The pair \((U_s,K_s)\) is the modal site generated by that declared occurrence. Two sites are equal only when their primitive rewrite constructors and occurrence addresses agree.
\end{definition}

\begin{definition}[Closure support for a site]
\label{def:closure-support}
A declared site records a displayed redex context; it does not say that
arbitrary rewrites may be propagated through that context.  A \emph{closure
support} for a site \(s\) is a chosen displayed contextual-closure constructor
\(\kappa_s\) in the operational presentation which transports a rewrite witness
in the hole through the displayed site context.  Schematically, for the same
ambient, index, and rely parameters as the site, it has the form
\[
  r:P\rw Q
  \quad\longmapsto\quad
  \kappa_s(r):K_s[P]\rw K_s[Q].
\]
After substituting ambient, index, and rely arguments, the same displayed constructor
therefore gives
\[
  \kappa_s(r):
  K_s[P][\vec c/\vec z,\vec a/\vec x,\vec b/\vec y]
  \rw
  K_s[Q][\vec c/\vec z,\vec a/\vec x,\vec b/\vec y].
\]
We write \(\kappa_s\Vdash K_s[-]\) when such a closure support has been
specified, and call \(s\) \emph{closure-supported}.  Closure support is separate
from site declaration: it is not inferred from a judgment \(P\ty\Dia C\), from
the fact that \(K_s[-]\) is a syntactic one-hole context, or from structural
equations.  It must be supplied by an explicit operational rule, such as a head
context rule for lambda calculus or a parallel-context rule for process
calculus.
\end{definition}

Because \(\mathcal{P}\) and \(\mathfrak{S}\) are finite, \(\Sites(\mathcal{P})\)
is finite.  The site specification also fixes the intended granularity of the
modal completion.  For the lambda beta rule, the product completion declares the
function-position occurrence and not the argument occurrence.  For the
asynchronous communication rule, the receive completion declares the input
occurrence and not necessarily the output occurrence.

The free variables of the selected subterm and its surrounding context determine
the modal interface, together with an ambient telescope:
\[
  \vec{x}_s
  =
  \FV(K_s)\cap\FV(U_s),
\]
\[
  \vec{y}_s
  =
  \FV(K_s)\setminus\FV(U_s),
\]
\[
  \vec{z}_s
  =
  \bigl(\FV(U_s)\setminus\FV(K_s)\bigr)
  \cup
  \bigl(\FV(R_i)\setminus\FV(L_i)\bigr).
\]
The variables \(\vec{x}_s\) are the \emph{index variables}; they occur both in
the term being typed and in the surrounding context, and they remain in the
external context of the modal type.  The variables \(\vec{y}_s\) are the
\emph{rely variables}; they are supplied by the surrounding context and
discharged by the modal type former.  The variables \(\vec{z}_s\) are the
\emph{ambient variables}; they are ordinary parameters of the selected subterm or
of the target which are not supplied by the one-hole context and are not
discharged by the modal type former.

For each site, let \(\Theta_s\) be the finite telescope of declarations for
\(\vec{z}_s\) inherited from the primitive rewrite context
\(\Gamma_i,\Delta_i\).  The ambient telescope is not part of the modal interface
and contributes no profile slots.  It is a side-context for the generated rules:
when the rules below write \(\Gamma,\Theta_s\), they mean an arbitrary surrounding
context \(\Gamma\) extended by declarations making all variables in \(\vec{z}_s\)
well formed.  If \(\vec{z}_s\) is empty, \(\Theta_s\) is omitted.

\subsection{Consistency relations on profiles}

For a site \(s\), write
\[
  I_s
  =
  \{i_x\mid x\in\vec{x}_s\}
  \amalg
  \{r_y\mid y\in\vec{y}_s\}
  \amalg
  \{\tau_s\}.
\]
The slot \(i_x\) records the sort chosen for the type of the index variable
\(x\), the slot \(r_y\) records the sort chosen for the type of the rely
variable \(y\), and \(\tau_s\) records the sort chosen for the target type
\(C\).  Ambient variables have no slots in \(I_s\): their declarations are
ordinary assumptions in \(\Theta_s\), not type-forming axes of the contextual
modality.  If there are \(m_s\) index variables and \(k_s\) rely variables, then
\(I_s\) has \(m_s+k_s+1\) elements.

We will use an explicit finite labeling of the type occurrences whose sorts are
controlled by these slots.  Let \(\TOcc_{\mathcal{P},s}\) be the finite set of
\emph{profile-relevant type occurrences} appearing in the displayed rule
templates generated from the site \(s\), together with the profile-relevant type
occurrences in displayed equations that are explicitly matched to such rule
occurrences.  There is a total slot-labeling map
\[
  \ell_s:\TOcc_{\mathcal{P},s}\longrightarrow I_s .
\]
For example, an occurrence of the type attached to an index or rely variable
\(v\) is labeled by the corresponding index or rely slot, and an occurrence of the
result type of the primitive rewrite target is labeled by the target slot
\(\tau_s\).  We also write
\[
  \mu_s:\TOcc_{\mathcal{P},s}\rightharpoonup \mathsf{TMeta}
\]
for the partial map that records the displayed type metavariable, when an
occurrence is literally named by one.

For each displayed equation \(E:L=R\) and site \(s\), let
\[
  \Match_s(E)\subseteq \TOcc_{\mathcal{P},s}\times \TOcc_{\mathcal{P},s}
\]
be the finite relation of profile-relevant type occurrences that the equation
schema matches for profile purposes. This relation is not the structural
congruence generated by \(E\); it is extra finite bookkeeping attached to the
displayed equation.  When no matching data are written, we use the following
syntactic default.  The default relation contains exactly those pairs \((o,o')\)
of occurrences on opposite sides of \(E\) such that either
\(\mu_s(o)=\mu_s(o')\) is defined, or \(o\) and \(o'\) occur at the same tree
address inside subterms with identical constructor heads and arities.  In both
cases the carriers must agree.  Structural equations that move occurrences to
different addresses, such as commutativity or associativity, therefore contribute
no such moved identifications unless their matching pairs are explicitly listed
in \(\Match_s(E)\).

The construction of the profile quotient is a finite syntactic preprocessing
step.  It uses the displayed finite presentation \(\mathcal{P}\), including the
finite relations \(\Match_s(E)\), not the full congruence or equational theory
generated by \(\mathcal{P}\).  Thus a listed equation contributes constraints
only through its displayed finite occurrence matching.  The same equation still
belongs to the equational layer and may be used for ordinary conversion in the
generated theory; that conversion role is separate from its profile-matching
role.  We do not enumerate consequences obtained by substitution into arbitrary
contexts, by transitivity of equality, or by rewriting modulo the listed
equations.

Each slot \(e\in I_s\) has a carrier, written \(\operatorname{car}(e)\):
\(\operatorname{car}(i_x)\) is the carrier of \(x\),
\(\operatorname{car}(r_y)\) is the carrier of \(y\), and
\(\operatorname{car}(\tau_s)=\Proc\).  The slot attached to a variable of carrier
\(X\) ranges over the two generated sorts \(\{\ast^X,\square^X\}\), and the
target slot ranges over \(\{\ast^\Proc,\square^\Proc\}\).  When all slots have
carrier \(\Proc\), we suppress the carrier superscripts and write the raw profile
set as
\[
  \{\ast,\square\}^{I_s}
  \cong
  \{\ast,\square\}^{m_s+k_s+1}.
\]
In general, the raw profile set is
\[
  \RawPrfl_{\mathcal{P}}(s)
  =
  \prod_{e\in I_s}\{\ast^{\operatorname{car}(e)},
                    \square^{\operatorname{car}(e)}\}.
\]

\begin{definition}[Profile-consistency algorithm]
For each site \(s\), form a finite undirected graph \(G_{\mathcal{P},s}\) with
vertex set \(I_s\), initialized with no edges.  The following finite clauses
propose edges.  A proposed edge is inserted exactly when it passes the carrier
check in clause \textup{(v)}.

\begin{enumerate}[label=\textup{(\roman*)}]
  \item \textbf{Shared type metavariable.} For occurrences
  \(o,o'\in\TOcc_{\mathcal{P},s}\), if \(\mu_s(o)\) and \(\mu_s(o')\) are defined
  and equal, propose the edge \(\ell_s(o)\sim\ell_s(o')\).  For example, if the
  same displayed type metavariable is used both as a binder annotation and as the
  domain of the generated modal type, the annotation slot and the modal-domain
  slot are identified.

  \item \textbf{Binder annotation.} For every Church-style refinement of a
  binding constructor, let \(o_A\) be the type occurrence of the added annotation
  and let \(o_x\) be the occurrence of the type of the variable bound by that
  constructor.  Propose the edge \(\ell_s(o_A)\sim\ell_s(o_x)\).  This clause is
  applied only to the finitely many binding-constructor occurrences appearing in
  the displayed source, selected subterm, surrounding context, and target of the
  primitive rewrite constructor under consideration.

  \item \textbf{Target variable.} Suppose the target \(R_i\) of the primitive
  rewrite constructor defining the site is a displayed occurrence of a variable
  \(v\in\vec{x}_s\cup\vec{y}_s\).  Let \(\sigma_v\) be the corresponding slot:
  \(\sigma_v=i_v\) if \(v\in\vec{x}_s\), and \(\sigma_v=r_v\) if
  \(v\in\vec{y}_s\).  Let \(\tau_s\) be the target slot.  Propose the edge
  \(\sigma_v\sim\tau_s\).  This records the fact that the
  target-typing premise for \(R_i\) uses the contextual declaration of the same
  variable.  The clause deliberately ranges only over index and rely variables:
  ambient variables in \(\Theta_s\) are external parameters and do not have
  profile slots.

  \item \textbf{Listed equation.} For each displayed equation
  \(E\) of \(\mathcal{P}\) and each matched pair
  \((o,o')\in\Match_s(E)\), propose the edge
  \(\ell_s(o)\sim\ell_s(o')\).  This is the only way listed equations constrain
  profile slots: the algorithm does not derive additional slot identifications
  from the equational theory generated by the presentation.

  \item \textbf{Carrier check.} A proposed edge \(e\sim e'\) is inserted only if
  \(\operatorname{car}(e)=\operatorname{car}(e')\), unless the presentation
  explicitly contains a displayed carrier identification making the comparison
  well sorted.  In the ordinary case, this means that \(\Proc\)-slots are
  compared only with \(\Proc\)-slots and \(\Nm\)-slots only with \(\Nm\)-slots.
\end{enumerate}

Let \(\Con_{\mathcal{P},s}\) be the reflexive, symmetric, transitive closure of
the inserted finite edge relation.  Equivalently, \(\Con_{\mathcal{P},s}\) is
computed by a union-find pass over the finite graph \(G_{\mathcal{P},s}\).
\end{definition}

The quotient profile space is the set of raw profiles that are constant on the
\(\Con_{\mathcal{P},s}\)-classes:
\[
  \Prfl_{\mathcal{P}}(s)
  =
  \left\{
    \omega\in\RawPrfl_{\mathcal{P}}(s)
    \;\middle|\;
    e\,\Con_{\mathcal{P},s}\,e'
    \Rightarrow
    \omega(e)=\omega(e')
  \right\}.
\]
Equivalently,
\[
  \Prfl_{\mathcal{P}}(s)
  \cong
  \prod_{[e]\in I_s/\Con_{\mathcal{P},s}}
  \{\ast^{\operatorname{car}(e)},\square^{\operatorname{car}(e)}\},
\]
where the carrier check makes \(\operatorname{car}(e)\) well defined on each
equivalence class.  Thus the algorithm for computing \(\Prfl_{\mathcal{P}}(s)\)
is:
\begin{enumerate}[label=\textup{(\arabic*)}]
  \item extract the finite site interface \(I_s\), the occurrence set
  \(\TOcc_{\mathcal{P},s}\), the slot-labeling map \(\ell_s\), and the displayed
  equation matchings \(\Match_s(E)\),
  \item add the finite set of carrier-correct edges proposed above,
  \item compute the equivalence classes of \(G_{\mathcal{P},s}\), and
  \item enumerate one independent \(\ast/\square\) choice for each equivalence
  class.
\end{enumerate}

\begin{definition}[Signature choice]
A signature choice for a site-marked finite presentation
\((\mathcal{P},\mathfrak{S})\) is a family of subsets
\[
  \Sig_s\subseteq\Prfl_{\mathcal{P}}(s)
  \qquad
  (s\in\Sites(\mathcal{P})).
\]
Equivalently, it is a subset
\[
  \Sig
  \subseteq
  \coprod_{s\in\Sites(\mathcal{P})}\Prfl_{\mathcal{P}}(s).
\]
\end{definition}

\begin{definition}[Based signature choice]
For each site \(s\), let
\[
  \omega^0_s
\]
denote the all-\(\ast\) profile, with the appropriate carrier superscripts.
A based signature choice is a signature choice such that
\[
  \omega^0_s\in\Sig_s
\]
for every declared site \(s\).
\end{definition}

\begin{remark}[Structural equations do not merge sites]
For the asynchronous communication source \(n!m\mid n?k\), the site
specification may name the left output occurrence, the right input occurrence,
both, or neither.  These names are occurrences in the displayed source tree.
The structural equation \(p\mid q=q\mid p\) identifies the two process terms
\(n!m\mid n?k\) and \(n?k\mid n!m\) in the equational layer, but it does
not identify the two declared sites, create a reversed-context site, or add a
profile edge.  Such effects occur only when the site specification or the
profile-matching data explicitly says so.
\end{remark}

Thus the input is not an arbitrary family of contexts and not an arbitrary family of sort tuples. The finite presentation together with its site specification determines the displayed sites, and the separate equation-matching data determines the quotient profile spaces. The user, or language designer, chooses which operational occurrences are sites, decides which structural equations, if any, should carry profile matchings, may annotate selected sites with explicit contextual-closure constructors for delayed elimination, and then chooses a subset of the quotient profile space at each declared site.

If \(\Prfl_{\mathcal{P}}(s)\) has \(N\) elements, then arbitrary signature choices form an \(N\)-dimensional Boolean hypercube.  Based signature choices form the face in which the all-\(\ast\) profile is present, hence have \(2^{N-1}\) vertices.  This based face is the convention used for the Lambda-cube and receive-hypercube examples below.

\begin{example}[Lambda profiles]
For the beta source
\[
  \app(\lam(t),y)\rw t(y),
\]
the product-site specification declares the function-position subterm \(U=\lam(t)\).  This is also the only occurrence selected by the default site policy, since the argument \(y\) is a variable occurrence and the whole source is excluded.  The declared site gives the context
\[
  K[-]=\app(\hole,y).
\]
There is one rely variable, \(y\), and no index variable. Thus the raw profile set is
\[
  \{\ast,\square\}^{2}.
\]
With no index slot, a profile is written \((\beta;\gamma)\), where \(\beta\)
is the sort of the argument type and \(\gamma\) is the sort of the result type.
If the presentation imposes no further consistency relation between the domain slot and target slot, the quotient profile space consists of
\[
  (\ast;\ast),
  \qquad
  (\ast;\square),
  \qquad
  (\square;\ast),
  \qquad
  (\square;\square).
\]
Choosing only \((\ast;\ast)\) gives the simply typed corner. Choosing all four gives the CoC corner for the product site.
\end{example}


\begin{example}[A nontrivial profile quotient]
Consider a toy operational presentation with constructors
\[
  \mathsf{mark}:\Proc\to\Proc,
  \qquad
  \mathsf{choose}:\Proc\times\Proc\times\Proc\to\Proc,
\]
and primitive rewrite constructor
\[
  x,y:\Proc
  \vdash
  \rho:
  \mathsf{choose}(\mathsf{mark}(x),y,x)\rw y.
\]
Declare the site
\[
  U=\mathsf{mark}(x),
  \qquad
  K[-]=\mathsf{choose}([-],y,x).
\]
Then \(x\) is an index variable, since it occurs both in \(U\) and in \(K[-]\),
and \(y\) is a rely variable, since it occurs in \(K[-]\) but not in \(U\).  The raw slots are
\[
  i_x,
  \qquad
  r_y,
  \qquad
  \tau,
\]
for the type of \(x\), the type of \(y\), and the target type.  Before imposing
consistency, the raw profiles are triples
\((\alpha,\beta;\gamma)\in\{\ast,\square\}^3\), where \(\alpha\) is the index sort and \(\beta\) is the rely sort.

The target of the rewrite is the displayed variable \(y\).  Clause
\textup{(iii)} of the profile-consistency algorithm therefore adds the edge
\(r_y\sim\tau\).  No displayed constraint mentions \(i_x\).  Hence the quotient
has two independent classes,
\[
  \{i_x\},
  \qquad
  \{r_y,\tau\},
\]
and the profile space is
\[
  \Prfl_{\mathcal{P}}(s)
  =
  \{(\alpha,\beta;\beta)
    \mid
    \alpha,\beta\in\{\ast,\square\}\}.
\]
Thus this site has four profiles rather than the eight profiles of the raw
three-slot cube.
\end{example}





\section{Generated presentations and structured categories}
\label{app:generated-categories}
This appendix records the full generated-presentation data summarized in Section~\ref{sec:hypercube-functor}.

We separate the generated typed presentation from the categorical object used
for the universal property.  If
\[
  A=(\mathcal{P},\mathfrak{S},\Sig)
\]
is a site-marked operational presentation with a signature choice, write
\(\mathcal{G}(A)\) for the finite typed presentation obtained by the
hypercube construction.  Outside categorical statements we also write
\(\Hyp(A)\) informally for this generated presentation; categorically,
\(\Hyp(A)\) will be the corresponding free structured object defined below.

\paragraph{The source category.}
Let \(\mathbf{PSig}\) be the category whose objects are triples
\[
  (\mathcal{P},\mathfrak{S},\Sig),
\]
where \(\mathcal{P}\) is a finite presentation of an operational theory,
\(\mathfrak{S}\) is a finite site specification, and \(\Sig\) is a
signature choice, i.e. a family of subsets
\[
  \Sig_s\subseteq \Prfl_{\mathcal{P}}(s)
  \qquad(s\in\Sites(\mathcal{P},\mathfrak{S})).
\]

A morphism
\[
  F:(\mathcal{P},\mathfrak{S},\Sig)
  \longrightarrow
  (\mathcal{P}',\mathfrak{S}',\Sig')
\]
in \(\mathbf{PSig}\) is a strict presentation morphism sending carrier
objects, constructors, relations, equations, profile-matching annotations,
closure-support annotations, and primitive rewrite constructors of
\(\mathcal{P}\) to corresponding displayed generators of
\(\mathcal{P}'\), preserving \(\Proc\), \(\Rew\), and \(s,t\).  Strictness
means that, for each primitive rewrite constructor \(\rho_i:L_i\rw R_i\),
\(F\) sends the displayed source tree \(L_i\) to the displayed source tree of
\(F(\rho_i)\) with occurrence addresses preserved, before quotienting by
structural equations.  Thus if \(o\in\mathfrak{S}_i\), then the image
occurrence \(F(o)\) lies in \(\mathfrak{S}'_{F(i)}\), and \(F\) induces
\[
  \Sites(F):\Sites(\mathcal{P},\mathfrak{S})
  \longrightarrow
  \Sites(\mathcal{P}',\mathfrak{S}').
\]
It also preserves matching data and closure support:
\[
  (o,o')\in\Match_s(E)
  \quad\Longrightarrow\quad
  (F(o),F(o'))\in\Match_{\Sites(F)(s)}(F(E)),
\]
and \(\kappa_s\Vdash K_s[-]\) is sent to
\(F(\kappa_s)\Vdash K_{\Sites(F)(s)}[-]\).  Hence \(F\) induces maps of
quotient profile spaces
\[
  \Prfl_{\mathcal{P}}(s)
  \longrightarrow
  \Prfl_{\mathcal{P}'}(\Sites(F)(s)).
\]
Finally, \(F\) is signature-preserving:
\[
  \omega\in\Sig_s
  \quad\Longrightarrow\quad
  F(\omega)\in\Sig'_{\Sites(F)(s)}.
\]
Identities and composition are inherited from strict presentation morphisms;
the displayed-address, matching, closure-support, and signature-preservation
conditions are stable under composition.

\paragraph{The generated typed presentation.}
On an object \(A=(\mathcal{P},\mathfrak{S},\Sig)\), the presentation
\(\mathcal{G}(A)\) is obtained by adjoining:
\begin{enumerate}[label=\textup{(\roman*)}]
  \item the two sort constants \(\ast^X,\square^X\) for each relevant carrier
  \(X\),
  \item the sort axiom \(\ast^X\ty^X\square^X\),
  \item proof-relevant typing evidence displays \(p_X:\Ev_X\to X\times X\),
  \item Church-style annotations for binding constructors,
  \item the one-step possibility type \(\Dia C\),
  \item for every site \(s\) with \(\Sig_s\neq\varnothing\), one contextual
modal type former associated to \(K_s[-]\),
  \item for every chosen quotient profile \(\omega\in\Sig_s\), the
corresponding profile-indexed formation, introduction, elimination,
source-subterm, and computation rule schemas,
  \item and no primitive delayed-\(\Dia\) rules.  Delayed elimination for closure-supported sites is an admissible metatheorem, not part of the finite generated presentation.
\end{enumerate}
On a morphism \(F:A\to A'\), the induced map \(\mathcal{G}(F)\) acts by
applying \(F\) to old syntax and rewrites.  It sends the contextual modality
associated to a site \(s\) to the contextual modality associated to
\(\Sites(F)(s)\).  For each chosen profile \(\omega\in\Sig_s\), it sends the
corresponding profile-indexed rule schema to the schema indexed by
\(F(\omega)\in\Sig'_{\Sites(F)(s)}\).  The functoriality of
\(\mathcal{G}\) is proved in Appendix~\ref{app:functoriality-adjunction}.

\paragraph{The target category.}
Let \(\mathbf{DOpPres}\) be the category of finite dependently typed
operational presentations and strict morphisms of such presentations.  A strict
morphism preserves the named carrier objects, constructors, relations,
equations, binding structure, typing judgments, and inference-rule schemas;
equivalently, it sends every rule of the source presentation to the
corresponding named rule, or to a specified derivation of that rule, in the
target presentation.

The target of the universal property is not the unstructured category
\(\mathbf{DOpPres}\), but the category of typed operational presentations
equipped with an interpretation of the generated hypercube structure.  Define
\(\mathbf{PDTOT}\) as follows.  An object is a triple
\[
  X=(A,\mathcal{T},\alpha_X)
\]
where \(A\in\mathbf{PSig}\), \(\mathcal{T}\in\mathbf{DOpPres}\), and
\[
  \alpha_X:\mathcal{G}(A)\longrightarrow\mathcal{T}
\]
is a strict morphism of typed presentations.  Thus \(\alpha_X\) chooses, in
\(\mathcal{T}\), interpretations of the generated sort constants, typing
evidence displays, Church annotations, possibility type, contextual modalities,
and rules belonging to \(A\).  The presentation \(\mathcal{T}\) may contain
additional generators, equations, or rules; freeness only concerns the named
structure selected by \(\alpha_X\).

A morphism
\[
  (F,h):(A,\mathcal{T},\alpha_X)
  \longrightarrow
  (B,\mathcal{T}',\alpha_Y)
\]
in \(\mathbf{PDTOT}\) consists of a morphism \(F:A\to B\) in
\(\mathbf{PSig}\) and a strict morphism
\[
  h:\mathcal{T}\to\mathcal{T}'
\]
in \(\mathbf{DOpPres}\) such that
\[
  h\circ\alpha_X
  =
  \alpha_Y\circ\mathcal{G}(F).
\]
Equivalently, the square
\[
\begin{array}{ccc}
  \mathcal{G}(A) & \xrightarrow{\mathcal{G}(F)} & \mathcal{G}(B) \\
  {\scriptstyle \alpha_X}\downarrow & & \downarrow{\scriptstyle \alpha_Y} \\
  \mathcal{T} & \xrightarrow{h} & \mathcal{T}'
\end{array}
\]
commutes.  Identities and composition are componentwise.  The forgetful
functor
\[
  U:\mathbf{PDTOT}\longrightarrow\mathbf{PSig}
\]
is defined by
\[
  U(A,\mathcal{T},\alpha_X)=A,
  \qquad
  U(F,h)=F.
\]

The hypercube functor is the free structured extension
\[
  \Hyp:\mathbf{PSig}\longrightarrow\mathbf{PDTOT},
  \qquad
  \Hyp(A)=(A,\mathcal{G}(A),\mathrm{id}_{\mathcal{G}(A)}),
\]
and on morphisms
\[
  \Hyp(F)=(F,\mathcal{G}(F)).
\]
Theorem~\ref{thm:free-hyp} states the adjunction in the main text, and Appendix~\ref{app:functoriality-adjunction} gives the full proof.

\begin{proposition}[Finiteness]
If \(\mathcal{P}\) is finite, \(\mathfrak{S}\) is a finite site specification, and \(\Sig\) is a signature choice for \((\mathcal{P},\mathfrak{S})\), then the underlying generated presentation \(\mathcal{G}(\mathcal{P},\mathfrak{S},\Sig)\), equivalently \(\Hyp(\mathcal{P},\mathfrak{S},\Sig)\) outside the categorical section, is finitely generated.
\end{proposition}

\begin{proof}
There are finitely many primitive rewrite constructors, each raw displayed source term has finitely many \(\Proc\)-subterm occurrences, and \(\mathfrak{S}\) chooses only finitely many of them. Hence \(\Sites(\mathcal{P})\) is finite. The free-variable lists \(\vec{x}_s\), \(\vec{y}_s\), and \(\vec{z}_s\), and hence the ambient telescope \(\Theta_s\), are finite for each site. Each displayed equation has only finitely many explicit or default profile matchings. For each site, the raw profile set \(\{\ast,\square\}^{m_s+k_s+1}\) is finite, and its consistency quotient is finite. Therefore the chosen subset \(\Sig_s\) is finite. The construction adds finitely many sort constants, typing evidence displays, Church-style refinements, modal formers, and rule schemas generated from this finite data.
\end{proof}



\section{Full proof of functoriality and adjunction}
\label{app:functoriality-adjunction}
This appendix supplies the details behind Theorem~\ref{thm:free-hyp}.

Let
\[
  F:A=(\mathcal{P},\mathfrak{S},\Sig)
  \longrightarrow
  A'=(\mathcal{P}',\mathfrak{S}',\Sig')
\]
be a morphism in \(\mathbf{PSig}\).  The generated-presentation map
\[
  \mathcal{G}(F):\mathcal{G}(A)\longrightarrow\mathcal{G}(A')
\]
is defined by structural recursion on generators and rules.  On old raw syntax,
primitive rewrites, equations, site declarations, profile matchings, and
closure-support annotations, \(\mathcal{G}(F)\) is \(F\).  On generated sorts
and typing evidence displays it sends
\[
  \ast^X\mapsto\ast^{F(X)},
  \qquad
  \square^X\mapsto\square^{F(X)},
  \qquad
  (M\ty^X A_0)\mapsto(F(M)\ty^{F(X)}F(A_0)).
\]
On the generated modality at site \(s\), it sends
\[
  \Mod{K_s}{\vec{x}\ty\vec{A}}{\vec{y}\ty\vec{B}}{C}
\]
to
\[
  \Mod{K_{\Sites(F)(s)}}{F(\vec{x})\ty F(\vec{A})}
       {F(\vec{y})\ty F(\vec{B})}{F(C)}.
\]
The rule schema indexed by \(\omega\in\Sig_s\) is sent to the rule schema
indexed by \(F(\omega)\in\Sig'_{\Sites(F)(s)}\).
This target former exists because \(F\) preserves declared sites and chosen
profiles.  The action on inference rules is obtained by applying \(F\) to every
premise and conclusion, including ambient telescopes and closure supports.  For example, the modal-introduction rule at \(s\) is sent to the
modal-introduction rule at \(\Sites(F)(s)\): the formation premises for
\(\vec A\), \(\vec B\), and \(C\), the premise forming \(P:\Proc\), the rewrite
premise \(r:K_s[P]\rw Q\), and the target-typing premise \(Q\ty C\) are all
mapped by applying \(F\) to the displayed terms, telescopes, and evidence
assumptions.  The conclusion is the image judgment
\[
  F(P)\ty
  \Mod{K_{\Sites(F)(s)}}{F(\vec{x})\ty F(\vec{A})}
       {F(\vec{y})\ty F(\vec{B})}{F(C)}.
\]
The same clause handles source-subterm introduction, modal elimination, and
computation equations.  Delayed elimination is not a generated rule to be mapped;
closure-support preservation instead ensures that the admissibility theorem for
source sites transports to the corresponding admissibility theorem for target sites.

\begin{proposition}[Functoriality of the generated presentation]
The assignment
\[
  \mathcal{G}:\mathbf{PSig}\longrightarrow\mathbf{DOpPres}
\]
is a functor.
\end{proposition}

\begin{proof}
For each \(F:A\to B\) in \(\mathbf{PSig}\), the preceding structural-recursion
clauses define a strict morphism \(\mathcal{G}(F):\mathcal{G}(A)\to
\mathcal{G}(B)\).  The only nontrivial well-definedness checks are the
generated components indexed by sites, profiles, ambient telescopes, and closure
supports.  These are precisely the preservation requirements in the definition
of \(\mathbf{PSig}\)-morphism: displayed addresses send sites to sites;
matching annotations send consistent profiles to consistent profiles; signature
preservation sends chosen profiles to chosen profiles; and closure-support
preservation sends the hypotheses needed for admissible delayed elimination to the
corresponding target hypotheses.  Therefore every generated
symbol, equation, and primitive rule of \(\mathcal{G}(A)\) has a named image in
\(\mathcal{G}(B)\).

If \(F=\mathrm{id}_A\), the recursive clauses send every old and generated
generator to itself, hence \(\mathcal{G}(F)=\mathrm{id}_{\mathcal{G}(A)}\).  If
\(F:A\to B\) and \(G:B\to C\), then \(\mathcal{G}(G)\circ\mathcal{G}(F)\) and
\(\mathcal{G}(G\circ F)\) agree on old generators by functoriality of strict
presentation morphisms, and they agree on generated generators because all
clauses are defined by applying the underlying map to the relevant carrier,
site, profile, telescope, and closure-support data.  Since \(\mathcal{G}(A)\)
is generated by exactly these old and generated symbols and rules, the two
morphisms are equal.
\end{proof}

\begin{proposition}[The structured target category]
The data in Appendix~\ref{app:generated-categories} define a category
\(\mathbf{PDTOT}\), and
\[
  U:\mathbf{PDTOT}\longrightarrow\mathbf{PSig}
\]
is a functor.
\end{proposition}

\begin{proof}
For an object \(X=(A,\mathcal{T},\alpha_X)\), the identity morphism is
\[
  (\mathrm{id}_A,\mathrm{id}_{\mathcal{T}}).
\]
The defining square commutes because
\[
  \mathrm{id}_{\mathcal{T}}\circ\alpha_X
  =
  \alpha_X\circ\mathcal{G}(\mathrm{id}_A).
\]
If
\[
  (F,h):X\to Y,
  \qquad
  (F',h'):Y\to Z,
\]
then their composite is \((F'\circ F,h'\circ h)\).  If
\(Y=(B,\mathcal{T}',\alpha_Y)\) and
\(Z=(C,\mathcal{T}'',\alpha_Z)\), the square for the composite commutes because
\[
  h'\circ h\circ\alpha_X
  =
  h'\circ\alpha_Y\circ\mathcal{G}(F)
  =
  \alpha_Z\circ\mathcal{G}(F')\circ\mathcal{G}(F)
  =
  \alpha_Z\circ\mathcal{G}(F'\circ F).
\]
Associativity and identity laws are inherited componentwise from
\(\mathbf{PSig}\) and \(\mathbf{DOpPres}\).  The functor \(U\) keeps only the
first component, so it preserves identities and composition.
\end{proof}

\begin{theorem}[Free hypercube completion]
\label{thm:free-hyp-full}
The functor
\[
  \Hyp:\mathbf{PSig}\longrightarrow\mathbf{PDTOT}
\]
is left adjoint to the forgetful functor
\[
  U:\mathbf{PDTOT}\longrightarrow\mathbf{PSig}.
\]
Equivalently, for every \(A\in\mathbf{PSig}\) and every
\(X=(B,\mathcal{T},\alpha_X)\in\mathbf{PDTOT}\), there is a natural bijection
\[
  \mathrm{Hom}_{\mathbf{PDTOT}}(\Hyp(A),X)
  \cong
  \mathrm{Hom}_{\mathbf{PSig}}(A,U(X)).
\]
\end{theorem}

\begin{proof}
First, \(\Hyp\) is a functor.  By definition,
\[
  \Hyp(A)=(A,\mathcal{G}(A),\mathrm{id}_{\mathcal{G}(A)}),
  \qquad
  \Hyp(F)=(F,\mathcal{G}(F)).
\]
The square required for \(\Hyp(F)\) to be a \(\mathbf{PDTOT}\)-morphism is
\[
\begin{array}{ccc}
  \mathcal{G}(A) & \xrightarrow{\mathcal{G}(F)} & \mathcal{G}(B) \\
  {\scriptstyle \mathrm{id}}\downarrow & & \downarrow{\scriptstyle \mathrm{id}} \\
  \mathcal{G}(A) & \xrightarrow{\mathcal{G}(F)} & \mathcal{G}(B),
\end{array}
\]
which commutes.  Identities and composition follow from the functoriality of
\(\mathcal{G}\).

Now fix \(A\in\mathbf{PSig}\) and
\(X=(B,\mathcal{T},\alpha_X)\in\mathbf{PDTOT}\).  A morphism
\[
  (F,h):\Hyp(A)\to X
\]
in \(\mathbf{PDTOT}\) consists of a morphism \(F:A\to B\) in
\(\mathbf{PSig}\) and a strict morphism \(h:\mathcal{G}(A)\to\mathcal{T}\) in
\(\mathbf{DOpPres}\) satisfying
\[
  h\circ \mathrm{id}_{\mathcal{G}(A)}
  =
  \alpha_X\circ\mathcal{G}(F).
\]
Thus \(h\) is forced to be \(\alpha_X\circ\mathcal{G}(F)\).  Consequently the
map
\[
  (F,h)\longmapsto F
\]
from \(\mathrm{Hom}_{\mathbf{PDTOT}}(\Hyp(A),X)\) to
\(\mathrm{Hom}_{\mathbf{PSig}}(A,U(X))\) is injective.  It is also surjective:
given \(F:A\to U(X)=B\), the pair
\[
  \bigl(F,\alpha_X\circ\mathcal{G}(F)\bigr)
\]
is a \(\mathbf{PDTOT}\)-morphism \(\Hyp(A)\to X\) because it satisfies the
commuting-square condition by construction.  Hence the displayed hom-set map is
a bijection.

Naturality is componentwise.  Precomposition by a morphism \(A'\to A\) in
\(\mathbf{PSig}\) replaces \(F\) by its composite on the left, and the
corresponding \(\mathbf{DOpPres}\)-component is forced to be the same composite
under \(\mathcal{G}\).  Postcomposition by a morphism
\((G,k):X\to Y\) in \(\mathbf{PDTOT}\) replaces \(F\) by \(G\circ F\); the
commuting-square condition for \((G,k)\) gives
\[
  k\circ\alpha_X\circ\mathcal{G}(F)
  =
  \alpha_Y\circ\mathcal{G}(G)\circ\mathcal{G}(F)
  =
  \alpha_Y\circ\mathcal{G}(G\circ F),
\]
which is exactly the forced second component for the composite.  Therefore the
bijection is natural in both variables.

The unit at \(A\) is the identity morphism
\[
  \eta_A=\mathrm{id}_A:A\longrightarrow U\Hyp(A)=A,
\]
and the counit at \(X=(B,\mathcal{T},\alpha_X)\) is
\[
  \epsilon_X=(\mathrm{id}_B,\alpha_X):
  \Hyp(U(X))=(B,\mathcal{G}(B),\mathrm{id}_{\mathcal{G}(B)})
  \longrightarrow
  (B,\mathcal{T},\alpha_X).
\]
The first triangle identity is
\[
  \epsilon_{\Hyp(A)}\circ\Hyp(\eta_A)
  =
  (\mathrm{id}_A,\mathrm{id}_{\mathcal{G}(A)}),
\]
and the second is
\[
  U(\epsilon_X)\circ\eta_{U(X)}
  =
  \mathrm{id}_{U(X)}.
\]
Therefore \(\Hyp\dashv U\).
\end{proof}

This theorem is the formal meaning of ``freely closes an operational theory
under the generated type-forming structure'': to interpret the free completion
in any structured typed operational theory \(X\), it is necessary and sufficient
to give a strict morphism of the underlying site-and-signature data into
\(U(X)\).  The remaining generated sorts, type formers, and rules are then
forced by the structure map \(\alpha_X\).





\section{Full proof-relevant evidence semantics and soundness proof}
\label{app:soundness}
This appendix supplies the proof details behind Theorem~\ref{thm:soundness} and
Proposition~\ref{prop:regular-may}.

The semantic statement has two layers.  The primary layer is proof-relevant:
typing judgments denote evidence objects, and modal judgments retain the rewrite
witnesses that justify them.  The optional second layer forgets that evidence
and forms proof-irrelevant may-predicates by taking images and dependent
universal operations.

Let \(\mathcal E\) be a category with the finite limits and specified
exponentials needed to interpret the underlying lambda-theory presentation.  A
proof-relevant model of the generated typed presentation \(\mathcal H(A)\) in
\(\mathcal E\) interprets each carrier \(X\) as an object, each constructor as a
morphism, and each equation as an equality of morphisms.  For each carrier
\(X\), the typing judgment is interpreted by a display
\[
  p_X:\Ev_X\longrightarrow X\times X,
  \qquad
  p_X=\langle\tm_X,\tp_X\rangle .
\]
The map \(p_X\) is not required to be monic.  The fiber over \((p,a)\) is the
object of evidence witnessing \(p\ty^X a\).

If a type term \(A:X\) is interpreted in context \(\Gamma\mid\Delta\) by a
morphism
\[
  \llbracket A\rrbracket_{\Gamma\mid\Delta}:
  \llbracket\Gamma\mid\Delta\rrbracket\longrightarrow X,
\]
then typed context extension is the pullback
\[
  \llbracket\Gamma,x:X\mid\Delta,e_x:x\ty^X A\rrbracket
  =
  (\llbracket\Gamma\mid\Delta\rrbracket\times X)
  \times_{X\times X}
  \Ev_X,
\]
where the comparison map
\(\llbracket\Gamma\mid\Delta\rrbracket\times X\to X\times X\) is
\(\langle\pi_2,\llbracket A\rrbracket_{\Gamma\mid\Delta}\pi_1\rangle\).
Thus a point of
\(\llbracket\Gamma,x:X\mid\Delta,e_x:x\ty^X A\rrbracket\) is a triple
\((c,x,e_x)\), where \(c\in\llbracket\Gamma\mid\Delta\rrbracket\) and
\(e_x\) is evidence that \(x\) has type \(A(c)\).  Weakening, exchange, and
substitution are interpreted by the usual pullback maps between these evidence
contexts.

A typing judgment
\[
  \Gamma\mid\Delta\vdash P\ty A
\]
is valid when there is an evidence morphism
\[
  e_P:\llbracket\Gamma\mid\Delta\rrbracket\longrightarrow\Ev_X
\]
satisfying
\[
  p_X e_P
  =
  \left\langle
    \llbracket P\rrbracket_{\Gamma\mid\Delta},
    \llbracket A\rrbracket_{\Gamma\mid\Delta}
  \right\rangle .
\]
For non-typing atomic judgments, such as rewrite judgments, validity is
interpreted by the corresponding evidence object supplied by the operational
presentation.  For example, a rewrite judgment \(r:P\rw Q\) is interpreted by a
morphism into \(\Rew\) whose source and target projections are
\(P\) and \(Q\).

\paragraph{The possibility display.}
For a type term \(C:\Proc\) in context \(\Gamma\mid\Delta\), the proof-relevant evidence
for \(P\ty\Dia C\) is the finite-limit object of tuples
\[
  (Q,r,e_C)
  \quad\text{with}\quad
  r:P\rw Q
  \quad\text{and}\quad
  e_C:Q\ty C.
\]
In the free generated model, this evidence display is defined to be the witness
object above.  Thus it has destructors to a reduct \(Q\), a rewrite witness
\(r:P\rw Q\), and target evidence \(e_C:Q\ty C\), together with the
corresponding introduction constructor.  In an arbitrary structured model, the
fiber of \(p_\Proc:\Ev_\Proc\to\Proc\times\Proc\) over \((P,\Dia C)\) is
required, by the chosen structure map, to carry the same introduction,
elimination, and computation structure.

The \(\Dia\)-introduction rule is interpreted by the tuple constructor
\[
  (Q,r,e_C)
  \longmapsto
  \operatorname{diaIntro}(Q,r,e_C)
  \in
  \Ev_\Proc(P,\Dia C).
\]
The \(\Dia\)-eliminator is interpreted by destructing an element of
\(\Ev_\Proc(P,\Dia C)\) into \(Q\), \(r\), and \(e_C\), and then applying the
interpretation of the derivation in the extended context.  The usual freshness
condition says that the final conclusion is pulled back from the smaller context
and does not mention the destructed variables.

\paragraph{Contextual modalities.}
For a declared site \(s\), write
\[
  M_s=\Mod{K_s}{\vec x\ty\vec A}{\vec y\ty\vec B}{C}.
\]
Evidence for \(P\ty M_s\) is proof-relevant evidence for the site interface.
The generated modal-elimination rule is interpreted by a morphism which, from
\[
  e_M:P\ty M_s
  \qquad\text{and}\qquad
  \vec e_b:\vec b\ty\vec B,
\]
produces evidence
\[
  \operatorname{elim}_s(e_M,\vec b,\vec e_b):
  K_s[P][\vec b/\vec y]
  \ty
  \Dia\bigl(C[\vec b/\vec y]\bigr).
\]
The generated modal-introduction rule supplies such an interface from open
one-step evidence in the rely context:
\[
  r:K_s[P]\rw Q,
  \qquad
  e_C:Q\ty C
  \quad\Longrightarrow\quad
  e_M:P\ty M_s.
\]
Semantically, this is the abstraction operation associated to the discharged
rely telescope, using the specified binding/exponential structure of the model,
or equivalently the named modal-introduction rule morphism of the structured
presentation.  No proof-irrelevant image is involved at this stage.

If \(s=(i,o)\) and the displayed source factorization is \(L_i=K_s[U_s]\), the
source-subterm rule is the special case in which the one-step witness is the
primitive rewrite constructor \(\rho_i\).  The computation equation says that
source-subterm introduction followed by modal elimination gives the same
\(\Dia\)-evidence as direct \(\Dia\)-introduction from \(\rho_i\) and the target
typing evidence.  Semantically this is an equality of morphisms into the
possibility evidence display.

\begin{theorem}[Soundness of the generated rules]
\label{thm:soundness-full}
Let \(A=(\mathcal P,\mathfrak S,\Sig)\) be an object of
\(\mathbf{PSig}\), and let \(\mathcal M\) be a structured proof-relevant model
of the generated typed presentation \(\mathcal H(A)\) in the sense above.  If a
judgment \(J\) is derivable from the generated rules of \(\mathcal H(A)\), then
\(\llbracket J\rrbracket\) is valid in \(\mathcal M\).  In particular, if
\(J\) is \(\Gamma\mid\Delta\vdash P\ty A\), then there is an evidence morphism
\[
  e_P:\llbracket\Gamma\mid\Delta\rrbracket\longrightarrow\Ev_X
\]
over \(\langle\llbracket P\rrbracket_{\Gamma\mid\Delta},
\llbracket A\rrbracket_{\Gamma\mid\Delta}\rangle\).
\end{theorem}

\begin{proof}
The proof is by induction on the derivation of \(J\).  The structural rules are
pullback facts.  Weakening is pullback along the projection from the extended
evidence context, and substitution is pullback along the section interpreting
the substituted term together with its typing evidence.  Equational conversion
is sound because equations of the presentation are interpreted as equal
morphisms, so evidence over one representative transports to evidence over the
other by the specified conversion rule.

For each non-structural generated rule, a structured model contains exactly the
morphism from the evidence object of its premises to the evidence object of its
conclusion specified by that rule schema.  The \(\Dia\)-rules use the rewrite
display \(\Rew\rightrightarrows\Proc\) and the target typing evidence.  The
contextual modal rules use the displayed factorization \(L_i=K_s[U_s]\), the
index-first telescope, and the ambient telescope \(\Theta_s\).  The
source-subterm rule is sound because the substituted primitive rewrite
constructor \(\rho_i\) is a point of the modal-introduction premise object.  The
induction hypotheses provide evidence for the premises; composing with the
corresponding rule morphism gives evidence for the conclusion.  Admissible
delayed elimination is handled separately in
Proposition~\ref{prop:soundness-admissible-delayed}.
\end{proof}

\begin{proposition}[Soundness of admissible delayed elimination]
\label{prop:soundness-admissible-delayed}
Let \(s\) be a site equipped with closure support \(\kappa_s\Vdash K_s[-]\).
For every \(n>0\), the admissible delayed-elimination principle of
Proposition~\ref{prop:admissible-delayed-modal-elim} is semantically sound in
every structured proof-relevant model of \(\mathcal H(A)\).
\end{proposition}

\begin{proof}
The proof follows the same induction on nested possibility evidence as the
syntactic admissibility proof.  A proof-relevant element of \(\Dia^n M\) is a
chain of \(n\) rewrite witnesses ending in evidence for \(M\).  The
closure-support constructor \(\kappa_s\) maps each witness in that chain to a
witness through the displayed context \(K_s[-]\).  After the transported chain
reaches a process of contextual modal type, ordinary modal elimination supplies
the final site-interaction witness.  The result is a chain of \(n+1\) witnesses
ending in evidence for \(C[\vec b/\vec y]\), which is exactly evidence for
\(\Dia^{n+1}(C[\vec b/\vec y])\).  In proof-irrelevant regular semantics this
becomes the corresponding composite of images.
\end{proof}

The preceding theorem is the soundness result needed for the free completion.
It does not require the semanticist to define \(\Dia C\) as a least
proof-irrelevant may-predicate, nor to define contextual modalities by
quantifying over all rely arguments after evidence has been forgotten.  Those
exact proof-irrelevant interpretations are additional semantic choices.

\begin{proposition}[Proof-irrelevant may is extra structure]
\label{prop:regular-may-full}
In a regular semantic category, the proof-irrelevant interpretation of the
one-step possibility predicate is obtained from the proof-relevant evidence
display by taking a regular image.  For contextual modalities with discharged
rely variables, an exact proof-irrelevant interpretation additionally requires
the corresponding dependent universal operation along the rely-variable
projection.
\end{proposition}

\begin{proof}
For an ordinary one-step possibility type, the proof-relevant evidence object in
context \(\Gamma\) is
\[
  \{(P,Q,r,e_C)\mid r:P\rw Q\text{ and }e_C:Q\ty C\}.
\]
Forgetting \(Q\), \(r\), and \(e_C\) gives a map to the object of processes
\(P\).  In a regular category its image is the proof-irrelevant predicate of
processes for which such evidence exists:
\[
  \operatorname{im}\left(
    \{(P,Q,r,e_C)\mid r:P\rw Q\text{ and }e_C:Q\ty C\}
    \longrightarrow
    \Proc
  \right).
\]
This proves the claim for \(\Dia\).

For a contextual modality with rely variables, the proof-irrelevant reading in
\(\mathbf{Set}\), at fixed index data \(\vec x\in\llbracket\vec A\rrbracket\),
has the form
\[
  p\in
  \llbracket\Mod{K_s}{\vec x\ty\vec A}{\vec y\ty\vec B}{C}\rrbracket
  \quad\Longleftrightarrow\quad
  \forall \vec y\in\llbracket\vec B(\vec x)\rrbracket.
  \exists Q,r,e_C.
  r:K_s[p][\vec y]\rw Q
  \text{ and }
  e_C:Q\ty C(\vec x,\vec y).
\]
The existential part is an image of rewrite-and-typing evidence.  The discharged
rely variables contribute the universal quantifier; categorically this is a
right adjoint to pullback along the rely-variable projection, or equivalent
dependent universal structure.  Thus the proof-relevant generated presentation
uses explicit evidence and its named rule morphisms, while exact
proof-irrelevant may semantics requires the additional logical structure just
described.
\end{proof}


\begin{thebibliography}{99}

\bibitem{AdamekRosickyVitale}
J.~Ad\'amek, J.~Rosick\'y, and E.~M.~Vitale.
\newblock \emph{Algebraic Theories: A Categorical Introduction to General Algebra}.
\newblock Cambridge University Press, 2011.

\bibitem{Barendregt1991}
H.~Barendregt.
\newblock Introduction to generalized type systems.
\newblock \emph{Journal of Functional Programming}, 1(2):125--154, 1991.

\bibitem{Barendregt1992Types}
H.~Barendregt.
\newblock Lambda calculi with types.
\newblock In S.~Abramsky, D.~M. Gabbay, and T.~S.~E. Maibaum, editors, \emph{Handbook of Logic in Computer Science}, volume~2, pages 117--309. Oxford University Press, 1992.

\bibitem{Boudol1992}
G.~Boudol,
\emph{Asynchrony and the Pi-calculus},
Research Report RR-1702, INRIA, 1992.

\bibitem{Cartmell1986}
J.~Cartmell.
\newblock Generalised algebraic theories and contextual categories.
\newblock \emph{Annals of Pure and Applied Logic}, 32:209--243, 1986.

\bibitem{FiorePlotkinTuri1999}
M.~Fiore, G.~Plotkin, and D.~Turi.
\newblock Abstract syntax and variable binding.
\newblock In \emph{Proceedings of the 14th Annual IEEE Symposium on Logic in Computer Science}, pages 193--202, 1999.

\bibitem{HondaTokoro1991}
K.~Honda and M.~Tokoro,
``An object calculus for asynchronous communication'',
in \emph{ECOOP '91: European Conference on Object-Oriented Programming},
Lecture Notes in Computer Science, vol.~512, Springer, 1991, pp.~133--147.

\bibitem{Jacobs1999}
B.~Jacobs.
\newblock \emph{Categorical Logic and Type Theory}.
\newblock North-Holland, 1999.

\bibitem{LambekScott1986}
J.~Lambek and P.~J.~Scott.
\newblock \emph{Introduction to Higher-Order Categorical Logic}.
\newblock Cambridge University Press, 1986.

\bibitem{Lawvere1963thesis}
F.~W. Lawvere,
\emph{Functorial Semantics of Algebraic Theories},
Ph.D. thesis, Columbia University, 1963.

\bibitem{LeiferMilner2000}
J.~J.~Leifer and R.~Milner.
\newblock Deriving bisimulation congruences for reactive systems.
\newblock In \emph{CONCUR 2000}, Lecture Notes in Computer Science 1877, pages 243--258. Springer, 2000.

\bibitem{MilnerParrowWalker1992}
R.~Milner, J.~Parrow, and D.~Walker,
``A calculus of mobile processes, I'',
\emph{Information and Computation} \textbf{100} (1992), no.~1, 1--40.

\bibitem{PlotkinSOS}
G.~D. Plotkin,
\emph{A structural approach to operational semantics},
The Journal of Logic and Algebraic Programming,
60--61:17--139, 2004.
doi: \url{https://doi.org/10.1016/j.jlap.2004.05.001}.

\bibitem{TrimbleMultisorted}
T. Trimble,
\emph{Multisorted Lawvere theories},
Todd Trimble's nLab pages, January 31, 2018.
Available at \url{https://ncatlab.org/toddtrimble/published/multisorted+Lawvere+theories}.

\bibitem{Whitehead1898}
A.~N. Whitehead,
\emph{A Treatise on Universal Algebra, with Applications. Vol. I},
Cambridge University Press, Cambridge, 1898.
\end{thebibliography}

\end{document}
